\documentclass[11pt,a4paper]{article}
\usepackage[margin=1in]{geometry}
\usepackage{amsmath,amssymb}
\usepackage{adjustbox}
\usepackage{fvextra}
\usepackage{pdflscape}
\usepackage[colorlinks=true,allcolors=blue]{hyperref}
\setlength{\parindent}{0pt}
\setlength{\parskip}{0.6em}
\emergencystretch=3em
\sloppy

\title{Link reflection positivity of the two-dimensional overlap, recovered asymptotically for the gauge-field marginal%
\thanks{Multi-DAC. Authors in alphabetical order. Clawd and Wren are AIs. How this work was done and checked is described in the README of
the repository \texttt{Multi-DAC/releases}, where the scripts for every computed figure are filed with this
paper.}}
\author{Clawd \and Clayton Iggulden-Schnell \and Wren}
\date{27 November 2026}

\begin{document}
\maketitle

\begin{abstract}
In two-dimensional U(1) lattice gauge theory on an $L\times L$ torus with $L$ even, with the Wilson plaquette action and one vector-like (Dirac) overlap fermion built on the domain-wall kernel $H_T$ at $a_5 = 1$ (antiperiodic in time), at fixed physical fermion mass and in a small physical box ($e\ell < 0.4727$ at bulk mass $m = 0.7$, where $e$ is the gauge coupling and $\ell$ the side of the box), we prove that for every sufficiently large $L$ the gauge-field marginal of the measure violates link reflection positivity by at most a constant times $L^{-3/4}\,\|F\|_\infty^2$, for every bounded function $F$ of the gauge links on one side of the reflection; the result is asymptotic in $L$, no explicit threshold is given, and the bound becomes non-trivial only at astronomically large $L$ (by Wren's estimate $L\sim 10^{138}$ to $10^{141}$, the upper end at the headline box). Part I (Wren) states the theorem and proves it with an auxiliary measure that is exactly link reflection positive and close in total variation to the true one; Part II (Clawd) is Lemma Z, the lower bound on the fermion determinant's average that Part I uses.
\end{abstract}


\section*{Part I. LANE 2 (Route (v)): Auxiliary Cut-Truncated DWF Measure and Asymptotic Non-Perturbative OS Recovery (\texttt{T1-OS'})}
\emph{Wren --- 2026-11-27 (Day 15, rev. 3 after Clawd's Letter 295 Section 8 refuter round and firsthand verification of Neuberger \texttt{hep-lat/\allowbreak{}9911004} \& Adams \texttt{hep-lat/\allowbreak{}9907005}). Companion to Clawd's \texttt{LANE1-Z-LEMMA.md} (Letters 290/89/295) and superseding \texttt{LANE2-TV-OS-PROOF.md} (Letter 90) after Clawd's cold audit (Letter 293).}
\emph{Status: AUDITED \& COMPLETE --- BOTH LANES REFEREED (Clawd's Letter 294 audited Route (v) with zero internal kills and all five wounds \texttt{(1')}, \texttt{(6')}, \texttt{(7)}, \texttt{(10)}, \texttt{(12)} folded in; Clawd's Letter 295 cold refuter on Lane 1 Section 8 confirmed the $O(L(\ln L)^{3/2})$ rate and fixed $L_0(H_T) = 32$ and Neuberger's $m=0.7$ ceiling to open $e\ell < 0.4727$; Neuberger \texttt{hep-lat/\allowbreak{}9911004} \& Adams \texttt{hep-lat/\allowbreak{}9907005} read firsthand, closing Row 11; ready for Release 11 \texttt{T1-OS'}).}


\subsection*{0. Idea \& Provenance (Why Route (v) Supersedes Route (iv) and Vitali)}
In \texttt{LANE2-TV-OS-PROOF.md} (Route (iv), Letter 90), we tried to prove $d_{\text{TV}}(\mu_{\text{ov}}^{(L)}, \mu_{\text{Wil}}^{(L)}) \to 0$ by comparing the overlap measure $\mu_{\text{ov}}^{(L)}$ to the matched massive Wilson measure $\mu_{\text{Wil}}^{(L)}$. Clawd's cold audit (Letter 293) killed Region 3 ($\exp(O((\ln L)^2))$ from Simon's $\det_4$ Lipschitz bound vs. $L^{-1/6}$) and the Section 6 $Q \ne 0$ fallback (via the $2 L^{C_0}$ ceiling), and our own follow-up caught the 2D Schwinger log-divergence of $\mathbb{E}_W[R_{\text{Wil}}]$ (Letter 92). Repairing Route (iv) via Vitali's theorem required proving an $L$-uniform $L^2$ moment ratio $\sup_L \mathbb{E}_W[R^2]/(\mathbb{E}_W R)^2 < \infty$ for both $R_{\text{ov}}$ and $R_{\text{Wil}}$, plus fixed-$q$ extensions on $Q \ne 0$.

{\bfseries\boldmath Route (v)} is directly inspired by the auxiliary-functional reflection-positivity mechanism in {\bfseries\boldmath Section 125 (\texttt{Lemma 125.1}--\texttt{125.2}, pp. 1007--1018)} of OpenAI's 15 November 2026 manuscript \emph{Continuum Yang--Mills Theory} (\texttt{yang-mills-group/\allowbreak{}OpenAI-informal-20261115}, Part VII, chiral $SU(5)$): instead of comparing $\mu_{\text{ov}}^{(L)}$ to a \emph{different} fermion action ($\mu_{\text{Wil}}^{(L)}$), compare $\mu_{\text{ov}}^{(L)}$ directly to its {\bfseries\boldmath own cut-truncated auxiliary DWF measure $\mu_{\text{aux}}^{(L)}$}:

\begin{itemize}
\item{} Keep the 5D Shamir DWF numerator $\det D_{\text{DW}}^{(L_5, m_f)}[U]$ {\bfseries\boldmath verbatim} (it is already \emph{exactly} link-reflection positive for every even $L$ and $L_5 \ge 1$!).
\item{} Truncate {\bfseries\boldmath only} the Pauli--Villars denominator $\ln \det D_{\text{PV}}^{(L_5)}[U]$ to a strictly $\Theta$-factorized sum $v_+[U_+] + v_-[U_-]$ of Dirichlet-ball traces on interior sites of $\Lambda_\pm$, paired with $\Theta$-factorized half-torus admissibility \& action caps $\mathbf{1}_{G_+^{\text{cap}}}[U_+] \mathbf{1}_{G_-^{\text{cap}}}[U_-]$.
\end{itemize}

Because the interior of $\Lambda_+ \cup \Lambda_-$ cancels \emph{identically} between $\mu_{\text{ov}}^{(L)}$ and $\mu_{\text{aux}}^{(L)}$ up to $O(L^{-9})$, the log-ratio $\ln W_{\text{aux}}[U] - \ln W_{\text{ov}}[U]$ on $G \cap G_{\text{act}}$ reduces to a {\bfseries\boldmath 1D cut-strip sum} $\Delta_{\text{cut}}[U]$ over only $O(L \ln L)$ plaquettes near the two reflection cuts! Consequently:

\begin{enumerate}
\item{} On the good strip event $E_{\text{strip}} = \{|\Delta_{\text{cut}}| \le L^{-3/4}\}$, $W_{\text{ov}}$ cancels {\bfseries\boldmath identically} between numerator and denominator --- giving $O(L^{-3/4})$ with {\bfseries\boldmath zero} Simon $\det_4$, {\bfseries\boldmath zero} Schatten-4, {\bfseries\boldmath zero} Borell--TIS, {\bfseries\boldmath zero} Gaussian RN, {\bfseries\boldmath zero} Schwinger log-divergence, and {\bfseries\boldmath zero} Vitali $L^2$ conjecture.
\item{} On the complement $(G \cap G_{\text{act}}) \setminus E_{\text{strip}}$, an {\bfseries\boldmath exact 2D plaquette-dependent character expansion} gives probability $\le \exp(-\Omega(L^{5/4}/\ln L))$, which crushes BOTH our polynomial ceiling $2 L^{C_0}$ AND Clawd's already-signed-off Lane 1 denominator $\mathbb{E}_W[W_{\text{ov}}] \ge \exp(-O(L(\ln L)^{3/2}))$!
\item{} Every step works {\bfseries\boldmath simultaneously on all topological sectors $Q \in \mathbb{Z}$ and all holonomies $\phi \in \mathbb{T}^2$} (zero Section 6 gap!).
\end{enumerate}

\emph{(Prior art note: Weingarten \& Challifour, ``Continuum limit of $\text{QED}_2$ on a lattice'', Ann. Phys. {\bfseries\boldmath 123} (1979), 61--101, and Weingarten, Ann. Phys. {\bfseries\boldmath 126} (1980), 154--175, constructed the continuum limit of 2D lattice QED with Wilson fermions; OpenAI's Nov 2026 manuscript introduces the auxiliary cut-local functional mechanism in 4D chiral $SU(5)$ via Balaban RG (unaudited, covering buffered even observables). What is proved here is an elementary, line-by-line checkable asymptotic total-variation bound for the 2D $H_T$ overlap gauge marginal at fixed physical box $e\ell < e\ell_0$ ($e\ell < 0.4727$ aligned with Lane 1's Neuberger ceiling, $e\ell < 0.4735$ with uniform gap at fixed $\varepsilon = 0.267$ via inline Lemma 2.0, and up to the open inline bound $e\ell < 0.5175$ via Lemma 2.0's $\varepsilon < 0.291819$), covering every bounded gauge observable because 2D super-renormalizability makes $\Delta_{\text{cut}} \sim \text{polylog}(L)/L \to 0$, using only 2D character expansion, Dirichlet Combes--Thomas localization, and Lane 1's Lemma Z.)}


\subsection*{1. Setup, Conventions, and Main Theorem}
Matching \texttt{LANE1-Z-LEMMA.md} Section 0 \& Section 8 (with our Letter 89 Section 8 fixes):

\begin{itemize}
\item{} {\bfseries\boldmath Lattice \& gauge measure:} Even $L \times L$ torus $\Lambda = \mathbb{Z}_L \times \mathbb{Z}_L$ ($L = 2T \ge 4$, $V = L^2$), links $U_\mu(t, x) = e^{i A_\mu(t, x)} \in U(1)$ ($\mu \in \{0, 1\}$, $t$ = Euclidean time, $x$ = space), principal plaquette angles $\theta_p \in (-\pi, \pi]$, plain Wilson measure $\mu_W(dU) \propto \prod_p e^{\beta \cos\theta_p} \prod_l dU_l$ at $\beta = V / (e\ell)^2$.
\item{} {\bfseries\boldmath Link reflection $\Theta$:} The two reflection cuts bisect the temporal links $U_0(L/2 - 1, x)$ and $U_0(L - 1, x)$, partitioning $\Lambda$ into $\Lambda_+ = \{t \in \{0, \dots, L/2 - 1\}\}$ and $\Lambda_- = \{t \in \{L/2, \dots, L - 1\}\}$. Link reflection acts on gauge configurations by

\begin{equation*}
\adjustbox{max width=\linewidth}{$\displaystyle (\theta U)_1(t, x) := U_1(L - 1 - t, x), \qquad (\theta U)_0(t, x) := U_0\big((L - 2 - t) \bmod L,\, x\big)^*,$}
\end{equation*}

and antilinearly on observables by $(\Theta F)(U) := \overline{F(\theta U)}$. Let $U_+$ denote all links with both endpoints in $\Lambda_+$ (i.e. $U_1(t, x)$ for $0 \le t \le L/2 - 1$ and $U_0(t, x)$ for $0 \le t \le L/2 - 2$), and $U_- := \theta(U_+)$ its reflected half; the crossing links are $U_{\text{cut}} = \{U_0(L/2 - 1, x), U_0(L - 1, x)\}$.
\item{} {\bfseries\boldmath Parameters:} Fix $0 < m < 1$ and $\varepsilon > 0$ such that

\begin{equation*}
\adjustbox{max width=\linewidth}{$\displaystyle 1 - m - \frac{\varepsilon}{4} > 0, \qquad \mu_0^2 := m^2 - \varepsilon - \frac{\varepsilon^2}{8} - \frac{\varepsilon^2}{2(1 - m - \varepsilon/4)} > 0.$}
\end{equation*}

\emph{(At our standard benchmark $m = 0.7$, our inline 2D Lemma 2.0 holds for every $0 < \varepsilon < \varepsilon_{\text{2D}}^* = 0.29181907\ldots$, strictly improving on Neuberger's $d$-dimensional bound $\sqrt{1 - (2+\sqrt{2})\varepsilon} - |1-m|$ (\texttt{hep-lat/\allowbreak{}9911004}, eq. at lines 592 \& 607, cited by Adams \texttt{hep-lat/\allowbreak{}9907005}), whose zero is $\varepsilon_{\text{Neu}}^* = m(2-m)/(2+\sqrt{2}) = 0.26653283\ldots$ In particular, at the concrete choice $\varepsilon = 0.267$, Neuberger's bound is negative ($-0.00267$), whereas inline Lemma 2.0 still gives $1 - m - \varepsilon/4 = 0.23325 > 0$ and a healthy uniform gap $\mu_0^2 = 0.061272 > 0$, $\mu_0 = 0.24753$.)}
Fix physical mass $\mu\ell > 0$ ($m_f := \mu\ell / L \in (0, 1)$ for $L > \mu\ell$) and physical box size

\begin{equation*}
\adjustbox{max width=\linewidth}{$\displaystyle 0 < e\ell < e\ell_0 := \frac{\varepsilon}{\sqrt{2 c_f(H_T)}},$}
\end{equation*}

where $c_f(H_T) = -\frac{1}{4\pi^2} \int_{\mathbb{T}^2} \ln|\lambda_0^{H_T}(k)|^2\, d^2k = 0.1589390$ at $m = 0.7$. Specifically, at $m = 0.7$:
\begin{itemize}
\item{} aligned with Lane 1 Section 8's Neuberger ceiling (Letter 295): open {\bfseries\boldmath $e\ell < 0.4727$} ($\varepsilon_{\text{Neu}}^*/\sqrt{2 c_f} = 0.472738$), on which inline Lemma 2.0 with fixed $\varepsilon = 0.267$ provides the uniform, non-degenerating gap $\mu_0 = 0.24753$ (in fact all the way to $e\ell < 0.4735$, since $0.267/\sqrt{2 c_f} = 0.473566$);
\item{} allowing $\varepsilon \in (e\ell\sqrt{2 c_f(H_T)}, \varepsilon_{\text{2D}}^*)$ via inline Lemma 2.0 alone (zero external citation): open {\bfseries\boldmath $e\ell < 0.5175$} ($\varepsilon_{\text{2D}}^*/\sqrt{2 c_f} = 0.517587$, with $L_0(H_T) = 32$ still valid in Lane 1 Section 8).
Let $G := \{\max_{p \in \Lambda} |1 - U_p| < \varepsilon\}$.
\end{itemize}

\item{} {\bfseries\boldmath Overlap operator ($H_T$ at $a_5 = 1$):} With Euclidean $\gamma_0 = \sigma_1, \gamma_1 = \sigma_2, \gamma_5 = \sigma_3$, antiperiodic temporal BC (seam $-1$ on $t = L - 1 \to 0$) and periodic spatial BC,

\begin{equation*}
\adjustbox{max width=\linewidth}{$\displaystyle D_W(M)[U] = (2 + M) I - \frac{1}{2}\sum_{\mu=0}^1 \big((I - \gamma_\mu) T_\mu + (I + \gamma_\mu) T_\mu^\dagger\big),$}
\end{equation*}


\begin{equation*}
\adjustbox{max width=\linewidth}{$\displaystyle H_W(-m)[U] = \gamma_5 D_W(-m)[U], \qquad H_T[U] = \gamma_5 D_W(-m)[U]\big(2 + D_W(-m)[U]\big)^{-1},$}
\end{equation*}


\begin{equation*}
\adjustbox{max width=\linewidth}{$\displaystyle D_{\text{ov}}^{H_T, (m_f)}[U] = (1 - m_f)\frac{1 + \gamma_5 \operatorname{sign}(H_T[U])}{2} + m_f.$}
\end{equation*}

Let $U_{\text{AP}}$ be the flat field with $A_0 \equiv 0$ (plus AP time seam) and $A_1 \equiv \pi / L$, and set

\begin{equation*}
\adjustbox{max width=\linewidth}{$\displaystyle W_{\text{ov}}[U] := \frac{\det D_{\text{ov}}^{H_T, (m_f)}[U]}{\det D_{\text{ov}}^{H_T, (m_f)}[U_{\text{AP}}]} > 0, \qquad \mu_{\text{ov}}^{(L)}(dU) := \frac{W_{\text{ov}}[U]}{\mathbb{E}_W[W_{\text{ov}}]}\, \mu_W(dU).$}
\end{equation*}

\end{itemize}

{\bfseries\boldmath Theorem (\texttt{T1-OS'}, Route (v) --- Gauge-Marginal Asymptotic OS Recovery).} \emph{Under the parameter hypotheses above (combining with Lane 1 Section 8's refereed lower bound $\mathbb{E}_W[W_{\text{ov}}] \ge \exp(-C_Z L (\ln L)^{3/2})$, which holds for every fixed $e\ell > 0$), for every sufficiently large even $L$ there exists an explicit probability measure $\mu_{\text{aux}}^{(L)}$ on $U(1)^{2V}$ (defined in Section 3.3) such that:}

\begin{enumerate}
\item{} \emph{(Exact finite-$L$ link reflection positivity of $\mu_{\text{aux}}^{(L)}$)} For every bounded measurable gauge observable $F(U_+)$ supported in $\Lambda_+$ (at any distance $d \ge 1$ from the reflection cuts, i.e., depending on any links in $U_+$),

\begin{equation*}
\adjustbox{max width=\linewidth}{$\displaystyle \mathbb{E}_{\text{aux}}^{(L)}\!\big[(\Theta F)\, F\big] \;\ge\; 0.$}
\end{equation*}

\item{} \emph{(Asymptotic total-variation convergence)}

\begin{equation*}
\adjustbox{max width=\linewidth}{$\displaystyle d_{\text{TV}}\!\left(\mu_{\text{ov}}^{(L)},\, \mu_{\text{aux}}^{(L)}\right) \;\le\; 2 L^{-3/4} + \exp\!\left(-c_*\, \frac{L^{5/4}}{\ln L}\right) \;=\; O\!\big(L^{-3/4}\big) \;\longrightarrow\; 0 \qquad (L \to \infty),$}
\end{equation*}

\emph{where $c_* := \frac{\mu_T}{144(e\ell)^2 C_{\text{PV}}(\mu_0)} > 0$ is given explicitly in Sections 3.2--3.3 and depends only on $(e\ell, m, \varepsilon)$.}
\item{} \emph{(Corollary --- non-perturbative asymptotic OS recovery for the gauge marginal)} For every bounded measurable gauge observable $F(U_+)$,

\begin{equation*}
\adjustbox{max width=\linewidth}{$\displaystyle \mathbb{E}_{\text{ov}}^{(L)}\!\big[(\Theta F)\, F\big] \;\ge\; -2\, d_{\text{TV}}\!\left(\mu_{\text{ov}}^{(L)},\, \mu_{\text{aux}}^{(L)}\right) \|F\|_\infty^2 \;=\; -O\!\big(L^{-3/4}\big) \|F\|_\infty^2 \;\longrightarrow\; 0 \qquad (L \to \infty).$}
\end{equation*}

\end{enumerate}

\emph{(Scope \& asymptoticity note: This theorem is an {\bfseries\boldmath asymptotic} statement in $L$ for the {\bfseries\boldmath gauge marginal} of the 2D vector-like $H_T$ overlap measure at fixed physical mass $\mu\ell > 0$ in a small physical box ($e\ell < 0.4727$ aligned with Lane 1's Neuberger ceiling, $e\ell < 0.4735$ at fixed $\varepsilon = 0.267$ with uniform gap $\mu_0 = 0.24753$, extending to open $e\ell < 0.5175$ via inline Lemma 2.0). It does not address fermionic correlators with inserted propagators $D_{\text{ov}}^{-1}$, nor the existence of the continuum limit of $\mu_{\text{ov}}^{(L)}$. Although every constant entering $c_*$ is written explicitly in Sections 3.2--3.3, at $\varepsilon = 0.267$ the elementary Combes--Thomas exponent $\gamma_0 \approx 1.10 \times 10^{-4}$ makes $C_{\text{PV}}(\mu_0) \approx 2.69 \times 10^{22}$, so the Piece B tail bound first drops below $1$ only at astronomical $L \sim 10^{138}\text{--}10^{141}$ --- far above simulable lattice sizes. Compared with OpenAI's Section 125 (4D chiral $SU(5)$, unaudited, buffered even observables), the present result is 2D, small enough to check line by line, and covers every bounded gauge observable without a buffer because 2D super-renormalizability makes the strip error $\Delta_{\text{cut}} \sim \operatorname{polylog}(L)/L \to 0$.)}


\subsection*{2. Lemma 1 --- Positivity, Exterior-Power Diamagnetic Bound, and Exact Link RP for 5D Shamir DWF}
Let $D_{\text{DW}}^{(L_5, m_f)}[U]$ be the standard 5D Shamir domain-wall operator at $a_5 = 1$ (matching \texttt{verify\_lane2\_route5\_bundle.py}) and $D_{\text{PV}}^{(L_5)}[U] := D_{\text{DW}}^{(L_5, 1)}[U]$ its Pauli--Villars partner, with the exact finite-$L_5$ Schur identity

\begin{equation*}
\adjustbox{max width=\linewidth}{$\displaystyle \frac{\det D_{\text{DW}}^{(L_5, m_f)}[U]}{\det D_{\text{PV}}^{(L_5)}[U]} = \det D_N^{(L_5, m_f)}[U], \qquad D_N^{(L_5, m_f)}[U] := (1 - m_f)\frac{1 + \gamma_5 \tanh(L_5 \tilde{H}_T[U] / 2)}{2} + m_f,$}
\end{equation*}

where $\tilde{H}_T := 2 \operatorname{artanh}(H_T)$ (valid whenever $|H_T| < 1$, which holds for every $U$ since $\operatorname{Re} D_W(2 - m) \succeq (2 - m) I$ and $H_T = \gamma_5(I - 2(2 + D_W(-m))^{-1})$; explicitly $|H_T|^2 = (D_W^\dagger D_W)((2 + D_W)^\dagger(2 + D_W))^{-1} < I$ because $(2 + D_W)^\dagger(2 + D_W) - D_W^\dagger D_W = 4 + 2(D_W + D_W^\dagger) \succeq 4(1 - m) I \succ 0$!).

{\bfseries\boldmath Lemma 1.} \emph{For every even $L \ge 2$, every $0 < m < 1$, $0 < m_f \le 1$, $L_5 \ge 1$, and every $U \in U(1)^{2V}$:}

\begin{enumerate}
\item{} \emph{(Strict positivity)} $\det D_{\text{DW}}^{(L_5, m_f)}[U] > 0$ and $\det D_{\text{PV}}^{(L_5)}[U] > 0$.
\item{} \emph{(Diamagnetic ceiling)} $\det D_{\text{DW}}^{(L_5, m_f)}[U] \le \det D_{\text{DW}}^{(L_5, m_f)}[U_{\text{AP}}]$.
\item{} \emph{(Exact link reflection positivity)} For any bounded real measurable function $h(U_+)$ depending only on $U_+$ and any bounded complex $F(U_+)$,

\begin{equation*}
\adjustbox{max width=\linewidth}{$\displaystyle \int_{U(1)^{2V}} e^{-S_W[U]} \det D_{\text{DW}}^{(L_5, m_f)}[U]\, h(U_+)\, h(\theta U_-)\, (\Theta F)(U)\, F(U)\, dU \;\ge\; 0.$}
\end{equation*}

\end{enumerate}

{\bfseries\boldmath Proof.}

\begin{enumerate}
\item{} {\bfseries\boldmath Strict positivity} (incorporating the 1-line repairs from Letter 293):
\begin{itemize}
\item{} First, for any 2D Wilson operator $D_W(M)[U]$ with $M > 0$: setting $X_\mu := \frac{1-\gamma_\mu}{2} T_\mu + \frac{1+\gamma_\mu}{2} T_\mu^\dagger$, the orthogonal projectors $P_\mu^\pm = \frac{1\pm\gamma_\mu}{2}$ make $X_\mu$ unitary, so $D_W(M) = (M + 2) I - (X_0 + X_1)$ has spectrum inside the disc $|z - (M + 2)| \le 2$. Hence every real eigenvalue is $\ge M > 0$, while non-real eigenvalues occur in complex conjugate pairs by $\gamma_5$-Hermiticity ($\gamma_5 D_W \gamma_5 = D_W^\dagger$), giving $\det D_W(M)[U] > 0$.
\item{} Second, in antiperiodic fifth-dimension Fourier modes $k_5 = \frac{2\pi(n + 1/2)}{L_5}$ ($n = 0, \dots, L_5 - 1$), $D_{\text{PV}}^{(L_5)}[U]$ block-diagonalizes into $\gamma_5 D_{\text{PV}}(k_5) = H_W(-m + u_5) + i\sin(k_5) I$ ($u_5 = 1 - \cos k_5$). Modes $k_5$ and $-k_5$ pair with $\det(D_{\text{PV}}(k_5) D_{\text{PV}}(-k_5)) = \det(H_W(-m + u_5)^2 + \sin^2 k_5) > 0$ ($\sin k_5 \ne 0$); if $L_5$ is odd, the single unpaired mode $k_5 = \pi$ has $u_5 = 2, \sin k_5 = 0$, giving $\det D_W(2 - m)[U] > 0$ (since $2 - m > 1 > 0$). Thus $\det D_{\text{PV}}^{(L_5)}[U] > 0$.
\item{} Third, $D_N^{(L_5, m_f)} = \frac{1+m_f}{2} I + \frac{1-m_f}{2} \gamma_5 \tanh(L_5 \tilde{H}_T/2)$ satisfies $\|\gamma_5 \tanh(L_5 \tilde{H}_T/2)\|_{\text{op}} < 1 \le \frac{1+m_f}{1-m_f}$ (for $0 < m_f \le 1$), so $\operatorname{Re} \operatorname{spec} D_N^{(L_5, m_f)} > 0$ and $\gamma_5$-Hermiticity gives $\det D_N^{(L_5, m_f)}[U] > 0$. Therefore $\det D_{\text{DW}}^{(L_5, m_f)} = \det D_{\text{PV}}^{(L_5)} \det D_N^{(L_5, m_f)} > 0$.
\end{itemize}

\item{} {\bfseries\boldmath Diamagnetic ceiling:} Rotate gamma matrices unitarily so $\gamma_0 = \sigma_3, \gamma_1 = \sigma_1, \gamma_5 = \sigma_2$. On slice $t$, the within-slice block is $\begin{pmatrix} C_t & A_t \\ -A_t^\dagger & C_t \end{pmatrix}$ with $C_t \succeq (1 - m) I_{L L_5} \succ 0$, giving the positive Hermitian $2 L L_5 \times 2 L L_5$ transfer matrix

\begin{equation*}
\adjustbox{max width=\linewidth}{$\displaystyle T_t = \begin{pmatrix} I & 0 \\ A_t^\dagger & I \end{pmatrix} \begin{pmatrix} C_t^{-1} & 0 \\ 0 & C_t \end{pmatrix} \begin{pmatrix} I & A_t \\ 0 & I \end{pmatrix} \succ 0,$}
\end{equation*}

and with unitary temporal gauge rotation $G_t = \operatorname{diag}(I_{L_5} \otimes \overline{U_0(t)}, I_{L_5} \otimes \overline{U_0(t)})$:

\begin{equation*}
\adjustbox{max width=\linewidth}{$\displaystyle \det D_{\text{DW}}^{(L_5, m_f)}[U] = \left(\prod_{t=0}^{L-1} \det C_t\right) \left(\prod_{t,x} U_0(t, x)\right)^{L_5} \operatorname{Tr}_{\Lambda^*}\!\left(\prod_{t=0}^{L-1} \Lambda^*(G_t) \Lambda^*(T_t)\right).$}
\end{equation*}

Since $\det D_{\text{DW}}^{(L_5, m_f)}[U] > 0$ by Part 1, taking moduli and applying non-commutative Schatten--H\"older ($|\operatorname{Tr}(\prod_{t=0}^{L-1} X_t)| \le \prod_t \|X_t\|_{S_L}$) plus unitary invariance of $\|\cdot\|_{S_L}$ gives

\begin{equation*}
\adjustbox{max width=\linewidth}{$\displaystyle \det D_{\text{DW}}^{(L_5, m_f)}[U] \;\le\; \prod_{t=0}^{L-1} \Big(\det D_{\text{DW}}^{(L_5, m_f)}\big[U^{\text{stat}}(U_t)\big]\Big)^{1/L}.$}
\end{equation*}

Gauge-transforming each static slice $U^{\text{stat}}(U_t)$ to constant spatial link $e^{i\Phi_t/L}$ and swapping $(x, t)$ gives $\det D_{\text{DW}}[U^{\text{stat}}(U_t)] = |\operatorname{Tr}_{\Lambda^*}(\Lambda^*(T_x)^L \Lambda^*(G_{\Phi_t}))| \le \operatorname{Tr}_{\Lambda^*}(\Lambda^*(T_x)^L) = \det D_{\text{DW}}[U_{\text{AP}}]$.
\item{} {\bfseries\boldmath Exact link reflection positivity:} Write the crossing links as $U_{\text{cut}} = \{U_0^{\text{cut}, 0}(x), U_0^{\text{cut}, 1}(x)\}$ on the two cuts ($t = L/2 - 1$ and $t = L - 1$). Let $M_+(U_+) := \Lambda^*(T_{L/2-1}) \Lambda^*(G_{L/2-2}) \cdots \Lambda^*(G_0) \Lambda^*(T_0)$ (with $\prod_{t \in \Lambda_+} (\det C_t) (\prod_x U_0(t, x))^{L_5}$ absorbed into the half-amplitude). Under link reflection $\theta$ ($A_1(L-1-t) = +A_1(t)$, $A_0(L-2-t) = -A_0(t)$), the lower half-transfer matrix satisfies $M_-(U_-) = M_+(\theta U_-)^\dagger$, while on each cut $c \in \{0, 1\}$, the diagonal unitary operator $\big(\prod_x U_0^{\text{cut}, c}(x)\big)^{L_5} \Lambda^*(G_{\text{cut}, c})$ acts on an exterior-algebra basis state $\alpha$ of local occupation numbers $q_{c, x}(\alpha) \in \{0, \dots, 2 L_5\}$ by the exact phase $\prod_x U_0^{\text{cut}, c}(x)^{L_5 - q_{c, x}(\alpha)}$. Meanwhile, each crossing plaquette on cut $c$ expands via the character series $e^{\beta \operatorname{Re} U_p} = \sum_{n_{c, x} \in \mathbb{Z}} I_{n_{c, x}}(\beta) \big(U_0^{\text{cut}, c}(x) U_0^{\text{cut}, c}(x+1)^*\big)^{n_{c, x}} U_1^+(c, x)^{n_{c, x}} \overline{U_1^-(c, x)}^{n_{c, x}}$ with $I_n(\beta) > 0$. Integrating over Haar measure $dU_{\text{cut}}$ projects onto the exact Gauss-law constraint $q_{c, x}(\alpha) - L_5 = n_{c, x} - n_{c, x-1}$ at every site $x$ on each cut $c \in \{0, 1\}$:

\begin{equation*}
\adjustbox{max width=\linewidth}{$\displaystyle \int dU_{\text{cut}}\, e^{-S_W[U]} \det D_{\text{DW}}^{(L_5, m_f)}[U] = \sum_{\mathbf{n}, \alpha, \beta} \delta_{\text{Gauss}}(\mathbf{n}, \alpha, \beta)\, \Big(\prod_{c, x} I_{n_{c, x}}(\beta)\Big)\, \overline{\Phi_{\mathbf{n}, \alpha, \beta}(\theta U_-)}\, \Phi_{\mathbf{n}, \alpha, \beta}(U_+),$}
\end{equation*}

where every coefficient $\prod_{c, x} I_{n_{c, x}}(\beta) > 0$ (verified to \texttt{2.3e-17} Hermiticity and strictly positive Gram spectrum on $6 \times 1$ and $4 \times 2$ tori in \texttt{verify\_lane2\_route5\_bundle.py}). Multiplying by $h(U_+) h(\theta U_-) (\Theta F)(U) F(U) = \overline{(h F)(\theta U_-)} (h F)(U_+)$ and integrating over $dU_+ dU_-$ yields $\sum_{\mathbf{n}, \alpha, \beta} \delta_{\text{Gauss}} (\prod I_{n_{c,x}}) \big|\int dU_+\, \Phi_{\mathbf{n}, \alpha, \beta}(U_+) h(U_+) F(U_+)\big|^2 \ge 0$. \emph{(Remark --- Clawd, Letter 294: when $h$ is gauge-invariant, as here, one may alternatively average over the gauge group on $\Lambda_+$ and $\Lambda_-$ separately; $\Theta$ intertwines the two projections, reducing any bounded $F(U_+)$ to gauge-invariant $P_+ F$, where $U_{\text{cut}}$ can be gauged to $1$ directly by transformations on $\Lambda_+$ sites alone.)} $\blacksquare$
\end{enumerate}


\subsection*{3. Lemma 2 --- Inline 5-Line Local Gap, Dirichlet $M_B$ Localization, and the Auxiliary Measure $\mu_{\text{aux}}^{(L)}$}
Recall $M_{k_5}[U] := H_W(-m + u_5)[U]^2 + \sin^2(k_5) I$ ($u_5 = 1 - \cos k_5 \in [0, 2]$), so $\det D_{\text{PV}}^{(L_5)}[U] = \prod_{k_5} (\det M_{k_5}[U])^{1/2}$.


\subsubsection*{3.1 Lemma 2.0 (Inline 5-line local Adams gap --- zero external black box)}
{\bfseries\boldmath Lemma 2.0.} \emph{Let $B \subseteq \Lambda$ be any subset of sites, let $B^{+2} := \{y \in \Lambda : d_\infty(y, B) \le 2\}$, and suppose $\max_{p \subset B^{+2}} |1 - U_p| \le \varepsilon$ (with $1 - m - \varepsilon/4 > 0$ and $\mu_0^2 := m^2 - \varepsilon - \frac{\varepsilon^2}{8} - \frac{\varepsilon^2}{2(1 - m - \varepsilon/4)} > 0$). Then for every spinor $\psi$ supported on $B$ and every $k_5 \in [-\pi, \pi]$:}

\begin{enumerate}
\item{} $\langle \psi, H_W(-m)[U]^2 \psi \rangle \ge \mu_0^2 \|\psi\|_2^2$,
\item{} $\langle \psi, M_{k_5}[U] \psi \rangle \ge c_{\text{PV}}^2 \|\psi\|_2^2$, where $c_{\text{PV}} := \mu_0 / 2 > 0$.
\emph{In particular, on $G$, $|H_W(-m)[U]| \ge \mu_0$, $|H_T[U]| \ge \mu_T := \mu_0 / (6 - m)$ (since $\|2 + D_W(-m)\|_{\text{op}} \le 6 - m$), and $M_{k_5}[U] \succeq c_{\text{PV}}^2 I$; and for any ball $B_r(x)$ with $B_{r+2}(x) \subset \Lambda_+$, the Dirichlet restriction $M_{k_5}^{(B)}[U_+] := P_{B_r(x)} M_{k_5}[U_+] P_{B_r(x)}$ depends only on links in $B_{r+1}(x) \subset \Lambda_+$ and satisfies $M_{k_5}^{(B)}[U_+] \succeq c_{\text{PV}}^2 I$ whenever $\max_{p \subset \Lambda_+} |1 - U_p| < \varepsilon$.}
\end{enumerate}

{\bfseries\boldmath Proof.} Write $W_\mu := \frac{1}{2}(T_\mu + T_\mu^\dagger)$, $S_\mu := \frac{1}{2i}(T_\mu - T_\mu^\dagger)$, $X_\mu := I - W_\mu = \frac{1}{2} \nabla_\mu^\dagger \nabla_\mu$ ($\nabla_\mu := I - T_\mu$). Then $D_W(-m) = (X_0 + X_1 - m) + i(\gamma_0 S_0 + \gamma_1 S_1)$, and squaring $H_W(-m)^2 = D_W(-m)^\dagger D_W(-m)$ using $X_\mu^2 + S_\mu^2 = 2 X_\mu$ gives the exact algebraic identity (verified to \texttt{4.8e-15} on random Haar $U$):

\begin{equation*}
\adjustbox{max width=\linewidth}{$\displaystyle H_W(-m)^2 = m^2 I + (1 - m)(\nabla_0^\dagger \nabla_0 + \nabla_1^\dagger \nabla_1) + \frac{1}{4}\big(\nabla_0^\dagger \nabla_1^\dagger \nabla_1 \nabla_0 + \nabla_1^\dagger \nabla_0^\dagger \nabla_0 \nabla_1\big) + E_{\text{scal}} + E_{\text{spin}},$}
\end{equation*}

where

\begin{equation*}
\adjustbox{max width=\linewidth}{$\displaystyle E_{\text{scal}} = \frac{1}{4}\Big(\nabla_0^\dagger[\nabla_0, \nabla_1^\dagger]\nabla_1 + \nabla_0^\dagger\nabla_1^\dagger[\nabla_0, \nabla_1] + \nabla_1^\dagger[\nabla_1, \nabla_0^\dagger]\nabla_0 + \nabla_1^\dagger\nabla_0^\dagger[\nabla_1, \nabla_0]\Big),$}
\end{equation*}


\begin{equation*}
\adjustbox{max width=\linewidth}{$\displaystyle E_{\text{spin}} = \frac{i}{2}\gamma_0\big(\nabla_1^\dagger[\nabla_1, S_0] + [\nabla_1^\dagger, S_0]\nabla_1\big) + \frac{i}{2}\gamma_1\big(\nabla_0^\dagger[\nabla_0, S_1] + [\nabla_0^\dagger, S_1]\nabla_0\big) + \gamma_0\gamma_1 [S_0, S_1].$}
\end{equation*}

Every commutator $[\nabla_0^{\pm}, \nabla_1^{\pm}] = [T_0^{\pm}, T_1^{\pm}]$ acting on $\psi$ or $\nabla_\mu \psi$ ($\operatorname{supp} \psi \subseteq B$) evaluates exclusively on plaquettes inside $B^{+2}$, where $|1 - U_p| \le \varepsilon$, so each of $[\nabla_0, \nabla_1]$, $[\nabla_0, \nabla_1^\dagger]$, $[\nabla_1, S_0]$, $[\nabla_0, S_1]$, and $[S_0, S_1]$ has operator norm $\le \varepsilon$ on those vectors. Applying Cauchy--Schwarz:

\begin{itemize}
\item{} $|\langle \psi, E_{\text{scal}} \psi \rangle| \le \frac{\varepsilon}{4}(\|\nabla_0\psi\|^2 + \|\nabla_1\psi\|^2) + \frac{1}{4}(\|\nabla_1\nabla_0\psi\|^2 + \|\nabla_0\nabla_1\psi\|^2) + \frac{\varepsilon^2}{8}\|\psi\|^2$,
\item{} $|\langle \psi, E_{\text{spin}} \psi \rangle| \le \varepsilon(\|\nabla_0\psi\| + \|\nabla_1\psi\|)\|\psi\| + \varepsilon\|\psi\|^2$.
\end{itemize}

Subtracting from the positive terms cancels $\frac{1}{4}(\|\nabla_1\nabla_0\psi\|^2 + \|\nabla_0\nabla_1\psi\|^2)$ identically; completing the square in $\|\nabla_\mu \psi\|$ with coefficient $a := 1 - m - \varepsilon/4 > 0$ ($\min_{u \ge 0} (a u^2 - \varepsilon u) = -\varepsilon^2 / (4a)$, summed over $\mu = 0, 1$) proves Part 1:

\begin{equation*}
\adjustbox{max width=\linewidth}{$\displaystyle \langle \psi, H_W(-m)^2 \psi \rangle \;\ge\; \left(m^2 - \varepsilon - \frac{\varepsilon^2}{8} - \frac{\varepsilon^2}{2(1 - m - \varepsilon/4)}\right) \|\psi\|^2 = \mu_0^2 \|\psi\|^2.$}
\end{equation*}

For Part 2 (every $\psi$ supported on $B$ and every $k_5 \in [-\pi, \pi]$):

\begin{itemize}
\item{} If $u_5 \le \mu_0/2$: $\|H_W(-m + u_5)\psi\| \ge \|H_W(-m)\psi\| - u_5\|\psi\| \ge (\mu_0/2)\|\psi\|$, so $\langle \psi, M_{k_5}\psi \rangle \ge (\mu_0^2/4)\|\psi\|^2$.
\item{} If $\mu_0/2 < u_5 \le 2 - \mu_0/2$: $\sin^2 k_5 = u_5(2 - u_5) \ge \mu_0^2/4$.
\item{} If $u_5 > 2 - \mu_0/2$: $M_+ := -m + u_5 > 1 - \mu_0/2 \ge 1/2$, and $\operatorname{Re} \langle \psi, D_W(M_+)\psi \rangle \ge M_+\|\psi\|^2 \implies \|H_W(M_+)\psi\| \ge \frac{1}{2}\|\psi\| \ge \frac{\mu_0}{2}\|\psi\|$. $\blacksquare$
\end{itemize}


\subsubsection*{3.2 Lemma 2a (Uniform Combes--Thomas, Dirichlet Schur Localization, and Local-Gauge Trace Bound)}
{\bfseries\boldmath Exact 9-block structure \& Uniform Combes--Thomas decay.} Expanding $H_W(M)[U]^2 = D_W(M)[U]^\dagger D_W(M)[U]$ with $D_W(M) = (2 + M) I - \sum_{\mu=0}^1 (P_\mu^- T_\mu + P_\mu^+ T_\mu^\dagger)$ ($M = -m + u_5 \in [-m, 2 - m]$, so $2 + M \in (1, 4)$), the orthogonal identities $P_\mu^+ P_\mu^- = P_\mu^- P_\mu^+ = 0$ annihilate all straight two-step hops $v \to v \pm 2 e_\mu$. Hence every row $v$ of $M_{k_5}[U] = H_W(-m + u_5)[U]^2 + \sin^2(k_5) I$ has at most {\bfseries\boldmath 9} nonzero $2 \times 2$ blocks:

\begin{itemize}
\item{} the on-site diagonal $(v, v)$: $\big((2+M)^2 + 4 + \sin^2 k_5\big) I_2$ (strictly independent of $U$, since $T_\mu T_\mu^\dagger = T_\mu^\dagger T_\mu = I$);
\item{} 4 nearest-neighbor blocks $|v - u|_1 = 1$: $-(2 + M) u_{vu} I_2$, of operator norm $2 + M < 4$;
\item{} 4 diagonal-corner blocks $u = v + s_0 e_0 + s_1 e_1$ ($s_0, s_1 \in \{\pm 1\}$, $|v - u|_1 = 2$): $P_0^{\pm} P_1^{\mp} (T_0^{s_0} T_1^{s_1})_{v, u} + P_1^{\pm} P_0^{\mp} (T_1^{s_1} T_0^{s_0})_{v, u}$, where $\|P_0^{\pm} P_1^{\mp}\|_{\text{op}} = 1/\sqrt{2} < 1$, giving block norm $\le \sqrt{2} < 2$.
\end{itemize}

In particular, for any rectangular ball $B = B_r(x)$, the Dirichlet restriction $M_{k_5}^{(B)}[U] := P_B M_{k_5}[U] P_B$ depends {\bfseries\boldmath exclusively} on links with both endpoints inside $B$ (every intermediate vertex of a corner path between $v, u \in B$ lies inside the rectangle $B$).

Now let $M$ be either $M_{k_5}[U]$ on $\Lambda$ ($U \in G$) or $M_{k_5}^{(B)}[U]$ on $B = B_r(x)$ ($\max_{p \subset B_{r+2}(x)} |1 - U_p| < \varepsilon$, $L > 4 r + 12$). For any fixed site $x$, let $d_x(y) := \min(|y_0 - x_0|_{\text{tor}}, r) + \min(|y_1 - x_1|_{\text{tor}}, r)$ (1-Lipschitz on $\Lambda$, and equal to $|y - x|_1$ identically on $B_r(x)$). For $\gamma \in [0, 1/2]$, the tilted operator $M_\gamma := e^{\gamma d_x} M e^{-\gamma d_x} = M + \delta_\gamma M$ has zero diagonal perturbation, nearest-neighbor perturbation $\le 4 (e^\gamma - 1) \le 4 e^{1/2} \gamma$ per neighbor, and corner perturbation $\le 2(e^{2\gamma} - 1) \le 4 e \gamma$ per corner. By Schur's test on the block row/column sums,

\begin{equation*}
\adjustbox{max width=\linewidth}{$\displaystyle \|\delta_\gamma M\|_{\text{op}} \;\le\; \big(16 e^{1/2} + 16 e\big)\gamma \;=:\; C_H \gamma \qquad \big(\gamma \in [0, 1/2],\; C_H \approx 69.872\big).$}
\end{equation*}

Setting $\gamma_0 := \min\!\big(1/2,\, c_{\text{PV}}^2 / (2 C_H)\big) > 0$, we obtain uniformly for all $t \ge 0$, $k_5 \in [-\pi, \pi]$, and $x, y$:

\begin{equation*}
\adjustbox{max width=\linewidth}{$\displaystyle \big\|(M + t)^{-1}_{x, y}\big\|_{\text{op}} \;\le\; \frac{2}{c_{\text{PV}}^2 + t} e^{-\gamma_0 d_x(y)}.$}
\end{equation*}


Fix the localization radius and fifth-dimension extent (for even $L$ large enough that $L > 4 r_0 + 12$):

\begin{equation*}
\adjustbox{max width=\linewidth}{$\displaystyle r_0 := \left\lceil \frac{7}{\gamma_0} \ln L \right\rceil, \qquad L_5 := \left\lceil \frac{5}{\mu_T} \ln L \right\rceil.$}
\end{equation*}


{\bfseries\boldmath Lemma 2a (Dirichlet Schur identity \& explicit local-gauge resolvent trace bound).} \emph{Define the explicit constant}

\begin{equation*}
\adjustbox{max width=\linewidth}{$\displaystyle C_3(\mu_0) := \frac{64 e^{\gamma_0}}{c_{\text{PV}}^2} + \frac{1824 e^{2\gamma_0}}{c_{\text{PV}}^4}$}
\end{equation*}

\emph{(independent of $L, L_5, k_5, Q, \phi$). Then for all even $L > 4 r_0 + 12$ and every $k_5 \in [-\pi, \pi]$:}

\begin{enumerate}
\item{} \emph{(Full-torus Dirichlet-to-torus localization)} For every $U \in G$ and {\bfseries\boldmath every site $x \in \Lambda$} (including cut-strip sites in $S_{\text{cut}}$), writing $B := B_{r_0}(x)$ and $M_{k_5}^{(B)}[U] := P_B M_{k_5}[U] P_B$,

\begin{equation*}
\adjustbox{max width=\linewidth}{$\displaystyle \Big|\operatorname{tr}_{2\times 2}\!\Big(\big(\ln M_{k_5}[U]\big)_{x, x} - \big(\ln M_{k_5}^{(B)}[U]\big)_{x, x}\Big)\Big| \;\le\; \frac{896 r_0}{c_{\text{PV}}^2} e^{-2\gamma_0 r_0} \;=\; O(L^{-13}).$}
\end{equation*}

\emph{Moreover, whenever $B_{r_0 + 2}(x) \subset \Lambda_+$ (resp. $\Lambda_-$), $M_{k_5}^{(B)}[U]$ depends exclusively on links in $U_+$ (resp. $U_-$), and its spectral gap and Combes--Thomas bound require only $\max_{p \subset \Lambda_+} |1 - U_p| < \varepsilon$ (resp. $\Lambda_-$).}
\item{} \emph{(Exact local-gauge quadratic trace bound)} Let $E_{\text{sp}}(B)$ denote the spatial links with both endpoints in $B = B_{r_0}(x)$, let $d(x, l) := \min(|x - p|_1, |x - q|_1)$ for $l = (p, q) \in E_{\text{sp}}(B)$, and let $a^{(x)}$ be the local axial-gauge field on the simply-connected rectangle $B$ anchored at $x$ ($a_0^{(x)} \equiv 0$, $a_1^{(x)}(t_x, y_1) \equiv 0$, $a_1^{(x)}(t, y_1) := \pm \sum_{s \text{ between } t_x, t} \theta_{s, y_1}$).
\begin{itemize}
\item{} For every $U$ with $\max_{p \subset B_{r_0 + 2}(x)} |1 - U_p| < \varepsilon$ (in particular, for every $U_+$ with $\max_{p \subset \Lambda_+} |1 - U_p| < \varepsilon$ and $B_{r_0 + 2}(x) \subset \Lambda_+$),

\begin{equation*}
\adjustbox{max width=\linewidth}{$\displaystyle \Big|\operatorname{tr}_{2\times 2}\!\Big(\big(\ln M_{k_5}^{(B)}[U]\big)_{x, x} - \big(\ln M_{k_5}^{(B)}[U_{\text{AP}}]\big)_{x, x}\Big)\Big| \;\le\; C_3(\mu_0) \sum_{l \in E_{\text{sp}}(B)} e^{-2\gamma_0 d(x, l)} \big(a_l^{(x)}\big)^2.$}
\end{equation*}

\item{} Consequently, for every $U \in G$ and {\bfseries\boldmath every site $x \in \Lambda$},

\begin{equation*}
\adjustbox{max width=\linewidth}{$\displaystyle \Big|\operatorname{tr}_{2\times 2}\!\Big(\big(\ln M_{k_5}[U]\big)_{x, x} - \big(\ln M_{k_5}[U_{\text{AP}}]\big)_{x, x}\Big)\Big| \;\le\; C_3(\mu_0) \sum_{l \in E_{\text{sp}}(B)} e^{-2\gamma_0 d(x, l)} \big(a_l^{(x)}\big)^2 + \frac{1792 r_0}{c_{\text{PV}}^2} e^{-2\gamma_0 r_0}.$}
\end{equation*}

\end{itemize}

\end{enumerate}

{\bfseries\boldmath Proof.}

\begin{enumerate}
\item{} Since $U \in G$ and $L > 4 r_0 + 12$, every $x \in \Lambda$ has $\max_{p \subset B_{r_0 + 2}(x)} |1 - U_p| < \varepsilon$, so Lemma 2.0 gives $M := M_{k_5}[U] \succeq c_{\text{PV}}^2 I$ and $M_{BB} := M_{k_5}^{(B)}[U] \succeq c_{\text{PV}}^2 I$. Write $\ln A = \int_0^\infty ((1 + t)^{-1} I - (A + t)^{-1})\, dt$ for both $M$ and $M_{BB}$; the $(1+t)^{-1} I_2$ cancels identically in the $(x, x)$ block. Blocking $M + t = \begin{pmatrix} M_{BB} + t & M_{B\bar B} \\ M_{\bar B B} & M_{\bar B \bar B} + t \end{pmatrix}$ gives the exact Schur identity

\begin{equation*}
\adjustbox{max width=\linewidth}{$\displaystyle (M + t)^{-1}_{BB} - (M_{BB} + t)^{-1} = -(M_{BB} + t)^{-1} M_{B\bar B} (M + t)^{-1}_{\bar B B}.$}
\end{equation*}

Since $M_{B\bar B}$ is supported on boundary sites $y \in \partial B$ ($\|y - x\|_\infty = r_0$, at most $8 r_0$ sites) coupled to $z \notin B$ ($\|z - x\|_\infty = r_0 + 1$), with $\sum_{z \notin B} \|M_{y, z}\|_{\text{op}} \le 2 \cdot 4 + 3 \cdot 2 = 14$ per site $y$ and $d_x(y) \ge r_0$, $d_x(z) \ge r_0$, Combes--Thomas on both $(M_{BB} + t)^{-1}_{x, y}$ and $(M + t)^{-1}_{z, x}$ yields

\begin{equation*}
\adjustbox{max width=\linewidth}{$\displaystyle \Big|\operatorname{tr}_{2\times 2}\!\Big(\big(\ln M\big)_{x, x} - \big(\ln M_{BB}\big)_{x, x}\Big)\Big| \;\le\; 2 \int_0^\infty \frac{4 \cdot (8 r_0 \cdot 14)}{(c_{\text{PV}}^2 + t)^2} e^{-2\gamma_0 r_0}\, dt = \frac{896 r_0}{c_{\text{PV}}^2} e^{-2\gamma_0 r_0} = O(L^{-13}).$}
\end{equation*}

\item{} For the Dirichlet-ball bound on $B = B_{r_0}(x)$ ($L > 4 r_0 + 12$ ensures $B$ does not wrap around the torus and is simply connected): gauge transformations act on $M_{k_5}^{(B)} = P_B M_{k_5} P_B$ by diagonal unitary conjugation $P_B D_g P_B$, leaving the $(x, x)$ block strictly invariant. Since $U^{(x)}$ has identical plaquette holonomies to $U$ on $B$ (any AP time seam inside $B$ has trivial plaquette curvature $(-1)(-1) = +1$ and is removed as pure gauge, while $U_{\text{AP}}|_B$ has $U_p \equiv 1$ and maps to $U^{(x)} \equiv 1$), we have $(\ln M_{k_5}^{(B)}[U])_{x, x} = (\ln M_1)_{x, x}$ and $(\ln M_{k_5}^{(B)}[U_{\text{AP}}])_{x, x} = (\ln M_0)_{x, x}$, where $M_0 := M_{k_5}^{(B)}[1]$ and $M_1 := M_{k_5}^{(B)}[U^{(x)}]$. Put $R_j(t) := (M_j + t)^{-1}$ and $\Delta M := M_1 - M_0 = \Delta^{(1)} M + \Delta^{(2)} M$ (splitting into the part linear in $a^{(x)}$ and the quadratic-and-higher remainder), and use the exact symmetric second-order resolvent identity

\begin{equation*}
\adjustbox{max width=\linewidth}{$\displaystyle R_0(t) - R_1(t) = R_0(t) \Delta M R_0(t) - R_0(t) \Delta M R_1(t) \Delta M R_0(t).$}
\end{equation*}

\begin{itemize}
\item{} {\bfseries\boldmath Site-by-site charge-conjugation vanishing of the linear trace} (incorporating Refuter Row 5's corrected identity, verified to \texttt{5.3e-16} for $M_{BB}$ and \texttt{1.6e-15} for $M$): with $\mathcal{C} = \sigma_2$, $\mathcal{C} D_W[U]^T \mathcal{C}^{-1} = D_W[U^*]$, hence $\mathcal{C} H_W[U]^T \mathcal{C}^{-1} = -D_W[U^*]\gamma_5 = -\gamma_5 H_W[U^*] \gamma_5$ and

\begin{equation*}
\adjustbox{max width=\linewidth}{$\displaystyle \mathcal{C} M[U]^T \mathcal{C}^{-1} = \gamma_5 M[U^*] \gamma_5$}
\end{equation*}

(and identically for $M_{BB} = P_B M P_B$ since $P_B$ is real, symmetric, and commutes with $\mathcal{C}, \gamma_5$). At $U \equiv 1$, $H_W[1]^2$ is a scalar multiple of $I_2$ in spinor space, so $M_0$ commutes with $\gamma_5$ and satisfies $\mathcal{C} M_0^T \mathcal{C}^{-1} = M_0$. Differentiating in $a^{(x)}$ at $0$ gives $\mathcal{C} (\Delta^{(1)} M)^T \mathcal{C}^{-1} = -\gamma_5 (\Delta^{(1)} M) \gamma_5$, hence $A := R_0(t) \Delta^{(1)} M R_0(t)$ obeys $\mathcal{C} A^T \mathcal{C}^{-1} = -\gamma_5 A \gamma_5$. Taking the $2 \times 2$ trace at $(x, x)$:

\begin{equation*}
\adjustbox{max width=\linewidth}{$\displaystyle \operatorname{tr}_{2\times 2}(A_{x, x}) = \operatorname{tr}_{2\times 2}(\mathcal{C} (A^T)_{x, x} \mathcal{C}^{-1}) = -\operatorname{tr}_{2\times 2}(\gamma_5 A_{x, x} \gamma_5) = -\operatorname{tr}_{2\times 2}(A_{x, x}) \implies \operatorname{tr}_{2\times 2}(A_{x, x}) \equiv 0$}
\end{equation*}

(verified to \texttt{2.06e-12} across all sites of $B$, including boundary sites!).
\item{} {\bfseries\boldmath Term 1 quadratic remainder ($\Delta^{(2)} M$):} In local axial gauge ($a_0 \equiv 0$), $(\Delta M)_{v, v} = 0$, nearest-neighbor spatial entries ($|v - u|_1 = 1$, link $l = (v, u)$) satisfy $\|(\Delta M)_{v, u}\|_{\text{op}} \le 4 |a_l|$ and $\|(\Delta^{(2)} M)_{v, u}\|_{\text{op}} \le 2 a_l^2$, and each diagonal-corner entry ($u = v + s_0 e_0 + s_1 e_1$) uses one spatial link $l_1$ incident to $v$ and one $l_2$ incident to $v + s_0 e_0$, with $\|(\Delta M)_{v, u}\|_{\text{op}} \le |a_{l_1}| + |a_{l_2}|$ and $\|(\Delta^{(2)} M)_{v, u}\|_{\text{op}} \le \frac{1}{2}(a_{l_1}^2 + a_{l_2}^2)$. Each spatial link $l = (p, q) \in E_{\text{sp}}(B)$ enters 2 nearest-neighbor ordered pairs (coeff $\le 2 a_l^2$, $|x-v|_1 + |x-u|_1 \ge 2 d(x, l)$) and 8 corner ordered pairs (coeff $\le \frac{1}{2} a_l^2$, $|x-v|_1 + |x-u|_1 \ge 2 d(x, l) - 1$). Applying Combes--Thomas on $R_0(t)$ and integrating $\int_0^\infty \frac{4\, dt}{(c_{\text{PV}}^2 + t)^2} = \frac{4}{c_{\text{PV}}^2}$ gives

\begin{equation*}
\adjustbox{max width=\linewidth}{$\displaystyle \int_0^\infty \Big|\operatorname{tr}_{2\times 2}\!\big((R_0(t) \Delta^{(2)} M R_0(t))_{x, x}\big)\Big|\, dt \;\le\; \frac{64 e^{\gamma_0}}{c_{\text{PV}}^2} \sum_{l \in E_{\text{sp}}(B)} e^{-2\gamma_0 d(x, l)} a_l^2.$}
\end{equation*}

\item{} {\bfseries\boldmath Term 2 (exact Hilbert-space bound --- single sum, exact $e^{-2\gamma_0 d}$ decay):} Let $E_x : \mathbb{C}^2 \to \ell^2(B; \mathbb{C}^2)$ be the site-$x$ inclusion and $\Psi(t) := \Delta M R_0(t) E_x$. Since $R_0(t)$ and $\Delta M$ are Hermitian and $0 \prec R_1(t) \preceq (c_{\text{PV}}^2 + t)^{-1} I$,

\begin{equation*}
\adjustbox{max width=\linewidth}{$\displaystyle 0 \;\le\; \operatorname{tr}_{2\times 2}\!\big((R_0(t) \Delta M R_1(t) \Delta M R_0(t))_{x, x}\big) = \operatorname{tr}_{2\times 2}\!\big(\Psi(t)^\dagger R_1(t) \Psi(t)\big) \;\le\; \frac{1}{c_{\text{PV}}^2 + t} \|\Psi(t)\|_{\text{HS}}^2 \;\le\; \frac{2}{c_{\text{PV}}^2 + t} \sum_{v \in B} \|\Psi(t)_v\|_{\text{op}}^2.$}
\end{equation*}

At each $v \in B$, $\Psi(t)_v = \sum_{u \in B} (\Delta M)_{v, u} R_0(t)_{u, x}$ involves at most 6 spatial links: the 2 spatial links $l$ incident to $v$ (each entering once at nearest neighbor with $|x - u|_1 \ge d(x, l)$ and twice at corners with $|x - u|_1 \ge d(x, l) - 1$, total weight $\le (4 + 2 e^{\gamma_0}) \frac{2}{c_{\text{PV}}^2 + t} e^{-\gamma_0 d(x, l)} |a_l| \le \frac{12 e^{\gamma_0}}{c_{\text{PV}}^2 + t} e^{-\gamma_0 d(x, l)} |a_l|$) and the 4 spatial links $l$ incident to $v \pm e_0$ (each entering once at a corner with $|x - u|_1 \ge d(x, l)$, weight $\le \frac{2}{c_{\text{PV}}^2 + t} e^{-\gamma_0 d(x, l)} |a_l|$). Applying $(\sum_{j=1}^6 b_j)^2 \le 6 \sum_{j=1}^6 b_j^2$ and summing over $v \in B$ (each link $l = (p, q)$ is incident to at most 2 sites $v \in B$ and to $v \pm e_0$ for at most 4 sites $v \in B$) gives

\begin{equation*}
\adjustbox{max width=\linewidth}{$\displaystyle \sum_{v \in B} \|\Psi(t)_v\|_{\text{op}}^2 \;\le\; \frac{6(2 \cdot 144 e^{2\gamma_0} + 4 \cdot 4)}{(c_{\text{PV}}^2 + t)^2} \sum_{l \in E_{\text{sp}}(B)} e^{-2\gamma_0 d(x, l)} a_l^2 \;\le\; \frac{1824 e^{2\gamma_0}}{(c_{\text{PV}}^2 + t)^2} \sum_{l \in E_{\text{sp}}(B)} e^{-2\gamma_0 d(x, l)} a_l^2.$}
\end{equation*}

Multiplying by $\frac{2}{c_{\text{PV}}^2 + t}$ and integrating $\int_0^\infty \frac{2\, dt}{(c_{\text{PV}}^2 + t)^3} = \frac{1}{c_{\text{PV}}^4}$ yields $\frac{1824 e^{2\gamma_0}}{c_{\text{PV}}^4} \sum_{l \in E_{\text{sp}}(B)} e^{-2\gamma_0 d(x, l)} a_l^2$. Adding Term 1 and Term 2 proves the Dirichlet-ball bound with $C_3(\mu_0) = \frac{64 e^{\gamma_0}}{c_{\text{PV}}^2} + \frac{1824 e^{2\gamma_0}}{c_{\text{PV}}^4}$. Finally, for every $U \in G$ and every $x \in \Lambda$, applying Part 1 to both $U$ and $U_{\text{AP}} \in G$ yields the full-torus bound with tail $2 \cdot \frac{896 r_0}{c_{\text{PV}}^2} e^{-2\gamma_0 r_0} = \frac{1792 r_0}{c_{\text{PV}}^2} e^{-2\gamma_0 r_0}$. $\blacksquare$
\end{itemize}

\end{enumerate}


\subsubsection*{3.3 Definition of the Auxiliary Measure $\mu_{\text{aux}}^{(L)}$ \& Corollary 2b}
Define the interior half-tori and cut strips:

\begin{equation*}
\adjustbox{max width=\linewidth}{$\displaystyle \Lambda_+^{\text{int}} := \big\{x = (t, x_1) \in \Lambda_+ : r_0 + 2 \le t \le L/2 - 1 - (r_0 + 2)\big\}, \qquad \Lambda_-^{\text{int}} := \theta(\Lambda_+^{\text{int}}),$}
\end{equation*}


\begin{equation*}
\adjustbox{max width=\linewidth}{$\displaystyle S_{\text{cut}} := \Lambda \setminus (\Lambda_+^{\text{int}} \cup \Lambda_-^{\text{int}}), \qquad S_{\text{cut}}^+ := \big\{p \in \Lambda : d_\infty(p, S_{\text{cut}}) \le r_0 + 2\big\}.$}
\end{equation*}

Note $|S_{\text{cut}}| \le 4 L (r_0 + 2)$ and $|S_{\text{cut}}^+| \le 2 L (4 r_0 + 10) = O(L \ln L)$.

Set the action-cap threshold $K := 4(c_f(H_T) + 1)$, define the {\bfseries\boldmath $\Theta$-factorized} half-torus capped-admissible sets

\begin{equation*}
\adjustbox{max width=\linewidth}{$\displaystyle G_+^{\text{cap}} := \left\{U_+ : \max_{p \subset \Lambda_+} |1 - U_p| < \varepsilon \;\;\text{and}\;\; \sum_{p \subset \Lambda_+} (1 - \cos\theta_p) \le K (e\ell)^2\right\}, \qquad G_-^{\text{cap}} := \theta(G_+^{\text{cap}}),$}
\end{equation*}

and the {\bfseries\boldmath $\Theta$-factorized} interior Pauli--Villars half-exponents (for $U_+ \in G_+^{\text{cap}}$, $U_- \in G_-^{\text{cap}}$):

\begin{equation*}
\adjustbox{max width=\linewidth}{$\displaystyle v_+[U_+] := \frac{1}{2} \sum_{k_5} \sum_{x \in \Lambda_+^{\text{int}}} \operatorname{tr}_{2\times 2}\!\Big(\big(\ln M_{k_5}^{(B_{r_0}(x))}[U_+]\big)_{x, x} - \big(\ln M_{k_5}^{(B_{r_0}(x))}[U_{\text{AP}, +}]\big)_{x, x}\Big), \qquad v_-[U_-] := v_+[\theta U_-].$}
\end{equation*}

\emph{(By exact $\theta$-covariance of $M_{k_5}^{(B)}$ --- verified to \texttt{9.3e-15} in \texttt{verify\_lane2\_route5\_bundle.py} --- $v_-[U_-]$ equals the exact same Dirichlet sum over $x \in \Lambda_-^{\text{int}}$.)}

Finally, define the {\bfseries\boldmath auxiliary weight} and {\bfseries\boldmath auxiliary probability measure}:

\begin{equation*}
\adjustbox{max width=\linewidth}{$\displaystyle W_{\text{aux}}[U] := \mathbf{1}_{G_+^{\text{cap}}}[U_+]\, \mathbf{1}_{G_-^{\text{cap}}}[U_-] \cdot \frac{\det D_{\text{DW}}^{(L_5, m_f)}[U]}{\det D_{\text{DW}}^{(L_5, m_f)}[U_{\text{AP}}]} \cdot \exp\!\big(-v_+[U_+] - v_-[U_-]\big), \qquad \mu_{\text{aux}}^{(L)}(dU) := \frac{W_{\text{aux}}[U]}{\mathbb{E}_W[W_{\text{aux}}]}\, \mu_W(dU).$}
\end{equation*}

\emph{(Strict positivity $\mathbb{E}_W[W_{\text{aux}}] > 0$ follows automatically from Theorem Part 2 below once $\mathbb{E}_W[|W_{\text{ov}} - W_{\text{aux}}|] < \mathbb{E}_W[W_{\text{ov}}]$, and is also immediate since $W_{\text{aux}} > 0$ on the open neighborhood of $U_{\text{AP}}$.)}

{\bfseries\boldmath Corollary 2b (Exact RP, global polynomial ceiling on $W_{\text{aux}}$, and cut-strip difference $\Delta_{\text{cut}}$).} \emph{Let $G_{\text{act}} := \{\sum_{p \in \Lambda} (1 - \cos\theta_p) \le K (e\ell)^2\}$, set $q := e^{-2\gamma_0} \in (0, 1)$, and define the explicit constants}

\begin{equation*}
\adjustbox{max width=\linewidth}{$\displaystyle C_{\text{geom}}(\gamma_0) := \frac{4 q (1 + q)}{(1 - q)^4}, \qquad C_{\text{PV}}(\mu_0) := C_3(\mu_0) C_{\text{geom}}(\gamma_0), \qquad K' := 3 K, \qquad C_0 := \frac{6 C_{\text{PV}}(\mu_0) K' (e\ell)^2}{\mu_T}.$}
\end{equation*}

\emph{Then for all sufficiently large even $L$:}

\begin{enumerate}
\item{} \emph{(Exact link RP)} $\mu_{\text{aux}}^{(L)}$ satisfies $\mathbb{E}_{\text{aux}}^{(L)}[(\Theta F) F] \ge 0$ for every bounded gauge observable $F(U_+)$ (immediate from Lemma 1(3) with $h(U_+) := \mathbf{1}_{G_+^{\text{cap}}}[U_+] e^{-v_+[U_+]}$!).
\item{} \emph{(Global polynomial ceiling on $W_{\text{aux}}$)} {\bfseries\boldmath Everywhere on $U(1)^{2V}$} (including on $\neg G$!),

\begin{equation*}
\adjustbox{max width=\linewidth}{$\displaystyle 0 \;\le\; W_{\text{aux}}[U] \;\le\; L^{C_0} \;\le\; 2 L^{C_0}.$}
\end{equation*}

\item{} \emph{(Polynomial ceiling on $W_{\text{ov}}$ on $G \cap G_{\text{act}}$)} Everywhere on $G \cap G_{\text{act}}$, $0 < W_{\text{ov}}[U] \le 2 L^{C_0}$.
\item{} \emph{(Cut-strip log-ratio on $G \cap G_{\text{act}}$)} Everywhere on $G \cap G_{\text{act}}$, $U \in G_+^{\text{cap}} \cap G_-^{\text{cap}}$ and

\begin{equation*}
\adjustbox{max width=\linewidth}{$\displaystyle \ln W_{\text{aux}}[U] - \ln W_{\text{ov}}[U] \;=\; \Delta_{\text{cut}}[U] + O(L^{-2}),$}
\end{equation*}

where

\begin{equation*}
\adjustbox{max width=\linewidth}{$\displaystyle \Delta_{\text{cut}}[U] := \frac{1}{2} \sum_{k_5} \sum_{x \in S_{\text{cut}}} \operatorname{tr}_{2\times 2}\!\Big(\big(\ln M_{k_5}[U]\big)_{x, x} - \big(\ln M_{k_5}[U_{\text{AP}}]\big)_{x, x}\Big)$}
\end{equation*}

satisfies the pointwise strip plaquette bound

\begin{equation*}
\adjustbox{max width=\linewidth}{$\displaystyle \big|\Delta_{\text{cut}}[U]\big| \;\le\; \frac{1}{2} L_5 C_{\text{PV}}(\mu_0) \sum_{p \in S_{\text{cut}}^+} \theta_p^2 + O(L^{-11}).$}
\end{equation*}

\end{enumerate}

{\bfseries\boldmath Proof.}

\begin{itemize}
\item{} For any subset of sites $X \subseteq \Lambda$, substitute $\big(a_1^{(x)}(t_l, y_l)\big)^2 \le |t_l - t_x| \sum_{s \text{ between } t_x, t_l} \theta_{s, y_l}^2$ into Lemma 2a(2) and interchange the sum over $x \in X, l \in E_{\text{sp}}(B_{r_0}(x))$ with the plaquette $p = (s_p, y_p)$. A fixed plaquette $p = (s_p, y_p)$ lies between $t_x$ and $t_l$ iff either $t_x = s_p - u$, $t_l = s_p + 1 + v$ or $t_x = s_p + 1 + u$, $t_l = s_p - v$ for integers $u, v \ge 0$ (2 cases, each with $|t_l - t_x| = u + v + 1$). Meanwhile $d(x, l) = (u + v + 1) + k$, where $k := \min(|x_1 - y_p|, |x_1 - y_p - 1|) \ge 0$ occurs for exactly 2 values of $x_1$ for each $k \ge 0$. Summing over $u, v, k \ge 0$ (putting $n = u + v + 1$, with $n$ pairs $(u, v)$ per $n$):

\begin{equation*}
\adjustbox{max width=\linewidth}{$\displaystyle \sum_{x \in X} \sum_{l \in E_{\text{sp}}(B_{r_0}(x)) \text{ covering } p} |t_l - t_x|\, e^{-2\gamma_0 d(x, l)} \;\le\; 4 \left(\sum_{k=0}^\infty q^k\right) \left(\sum_{n=1}^\infty n^2 q^n\right) = \frac{4 q (1 + q)}{(1 - q)^4} = C_{\text{geom}}(\gamma_0).$}
\end{equation*}

On $G_+^{\text{cap}}$, $|\theta_p| \le \theta_\varepsilon < \pi/3 \implies \theta_p^2 \le 3(1 - \cos\theta_p)$, so $\sum_{p \subset \Lambda_+} \theta_p^2 \le 3 K (e\ell)^2 = K'(e\ell)^2$, giving $|v_+[U_+]| \le \frac{1}{2} L_5 C_{\text{PV}}(\mu_0) K'(e\ell)^2$. Adding the identical bound on $G_-^{\text{cap}}$ and using $L_5 = \lceil \frac{5}{\mu_T} \ln L \rceil \le \frac{6}{\mu_T} \ln L$ (valid for all $L \ge e^{\mu_T}$, hence all $L \ge 2$) gives $|v_+[U_+] + v_-[U_-]| \le L_5 C_{\text{PV}}(\mu_0) K'(e\ell)^2 \le C_0 \ln L$. Multiplying by $\det D_{\text{DW}}^{(L_5)}[U] / \det D_{\text{DW}}^{(L_5)}[U_{\text{AP}}] \le 1$ (Lemma 1(2)) yields $0 \le W_{\text{aux}}[U] \le L^{C_0}$ {\bfseries\boldmath everywhere on $U(1)^{2V}$}.
\item{} On $G$, $|H_T[U]| \ge \mu_T \implies |\tilde{H}_T[U]| \ge 2\mu_T \implies \|D_{\text{ov}}^{H_T}[U] - D_N^{(L_5, m_f)}[U]\|_{\text{op}} \le e^{-L_5 \mu_T} \le L^{-5}$. Since every singular value of $D_{\text{ov}}^{H_T}$ is $\ge m_f = \mu\ell/L$, $|\ln\det D_{\text{ov}}^{H_T}[U] - \ln\det D_N^{(L_5, m_f)}[U]| \le 2 V \frac{L}{\mu\ell} L^{-5} = O(L^{-2})$ uniformly on $G$ (and at $U_{\text{AP}}$). Combining with Lemma 1(2) and summing Lemma 2a(2) over all $x \in \Lambda$ gives Part 3 ($W_{\text{ov}}[U] \le 2 L^{C_0}$ on $G \cap G_{\text{act}}$).
\item{} For Part 4: on $G \cap G_{\text{act}}$, subtracting $\ln W_{\text{ov}} = \ln(\det D_{\text{DW}} / \det D_{\text{DW}}^{\text{AP}}) - \ln(\det D_{\text{PV}} / \det D_{\text{PV}}^{\text{AP}}) + O(L^{-2})$ from $\ln W_{\text{aux}}$ cancels the $\det D_{\text{DW}}$ ratio identically; every interior site $x \in \Lambda_+^{\text{int}} \cup \Lambda_-^{\text{int}}$ cancels up to $O(L^{-13})$ by Lemma 2a(1) (total over $V$ sites and $L_5$ modes: $O(L^{-11})$), leaving $\Delta_{\text{cut}}[U] + O(L^{-2})$. Summing Lemma 2a(2) over $x \in S_{\text{cut}}$ and $k_5$ yields the plaquette bound on $S_{\text{cut}}^+$ (with tail $\frac{1}{2} L_5 |S_{\text{cut}}| \cdot O(L^{-13}) = O(L^{-11})$). $\blacksquare$
\end{itemize}


\subsection*{4. Lemma 3 --- Exact Plaquette-Dependent 2D Character Expansion \& Strip Tail Bound}
{\bfseries\boldmath Lemma 3.0 (1-line Vandermonde--Tur\'an Bessel decay).} \emph{For all $x > 0$ and all integers $n \ge 1$, $I_n(x)^2 > I_{n-1}(x) I_{n+1}(x)$, and for all $n \ge 2$,}

\begin{equation*}
\adjustbox{max width=\linewidth}{$\displaystyle \frac{I_n(x)}{I_0(x)} \;\le\; \exp\!\left(-\frac{(n - 1)^2}{2(x + n)}\right).$}
\end{equation*}


{\bfseries\boldmath Proof.} In the power series $I_n(x) = \sum_{j=0}^\infty \frac{(x/2)^{n+2j}}{j!(n+j)!}$, the coefficient of $(x/2)^{2n+2m}$ in $I_n(x)^2 - I_{n-1}(x) I_{n+1}(x)$ is, by Vandermonde's identity $\sum_{j=0}^m \binom{m}{j}\binom{2n+m}{n+j} = \binom{2n+2m}{n+m}$,

\begin{equation*}
\adjustbox{max width=\linewidth}{$\displaystyle \frac{1}{m!(2n+m)!}\left[\binom{2n+2m}{n+m} - \binom{2n+2m}{n+m-1}\right] = \frac{1}{m!(2n+m)!(n+m+1)}\binom{2n+2m}{n+m} > 0,$}
\end{equation*}

proving $I_n(x)^2 > I_{n-1}(x) I_{n+1}(x)$ term-by-term! Thus $r_k := I_k(x) / I_{k-1}(x) \in (0, 1)$ is strictly decreasing in $k \ge 1$. Substituting $r_{k-1} > r_k$ into the exact Bessel recurrence $\frac{1}{r_{k-1}} - r_k = \frac{2(k-1)}{x}$ ($k \ge 2$) yields $\frac{1}{r_k} - r_k > \frac{2(k-1)}{x}$, so

\begin{equation*}
\adjustbox{max width=\linewidth}{$\displaystyle r_k < \frac{x}{(k - 1) + \sqrt{(k - 1)^2 + x^2}} = \exp\!\left(-\operatorname{arsinh}\frac{k - 1}{x}\right) \le \exp\!\left(-\frac{k - 1}{x + k - 1}\right) \qquad (k \ge 2)$}
\end{equation*}

(since $\operatorname{arsinh}(u) = \int_0^u (1 + s^2)^{-1/2}\, ds \ge \frac{u}{\sqrt{1 + u^2}} \ge \frac{u}{1 + u}$ for all $u \ge 0$). Multiplying over $k = 2, \dots, n$ (with $r_1 < 1$) and using $\sum_{k=2}^n (k - 1) = \frac{n(n-1)}{2} \ge \frac{(n-1)^2}{2}$ gives $I_n(x)/I_0(x) \le \exp(-\frac{(n-1)^2}{2(x+n)})$. $\blacksquare$

{\bfseries\boldmath Lemma 3 (Exact global action cap \& cut-strip action cap).} \emph{For all sufficiently large $L$:}

\begin{enumerate}
\item{} \emph{(Global action cap --- Lemma 3a)} With $K = 4(c_f(H_T) + 1)$, $\mu_W(\neg G_{\text{act}}) \le e^{-(c_f(H_T) + 1) V}$.
\item{} \emph{(Cut-strip tail)} Let $E_{\text{strip}} := \{|\Delta_{\text{cut}}[U]| \le L^{-3/4}\}$, and set the explicit constant $c_* := \frac{\mu_T}{144(e\ell)^2 C_{\text{PV}}(\mu_0)} > 0$. Then for all sufficiently large $L$,

\begin{equation*}
\adjustbox{max width=\linewidth}{$\displaystyle \mu_W\big((G \cap G_{\text{act}}) \setminus E_{\text{strip}}\big) \;\le\; \exp\!\left(-2 c_*\, \frac{L^{5/4}}{\ln L}\right).$}
\end{equation*}

\end{enumerate}

{\bfseries\boldmath Proof.} For any subset of plaquettes $S \subseteq \Lambda$ and any $\alpha \in (0, 1)$, expanding $e^{\beta_p \cos\theta_p} = \sum_{n_p \in \mathbb{Z}} I_{n_p}(\beta_p) e^{i n_p \theta_p}$ with $\beta_p = (1 - \alpha)\beta$ on $S$ and $\beta_p = \beta$ on $\Lambda \setminus S$ and integrating every link $U_l$ forces $n_p \equiv n$ across the 2D torus, yielding the {\bfseries\boldmath exact} identity (verified vs. Monte Carlo in \texttt{verify\_lane2\_route5\_bundle.py}):

\begin{equation*}
\adjustbox{max width=\linewidth}{$\displaystyle \mathbb{E}_W\!\left[\exp\!\Big(\alpha \beta \sum_{p \in S} (1 - \cos\theta_p)\Big)\right] = \frac{\sum_{n \in \mathbb{Z}} \operatorname{ive}\big(n, (1-\alpha)\beta\big)^{|S|} \operatorname{ive}(n, \beta)^{V - |S|}}{\sum_{n \in \mathbb{Z}} \operatorname{ive}(n, \beta)^V},$}
\end{equation*}

where $\operatorname{ive}(n, x) := e^{-x} I_n(x) > 0$. Take $\alpha = 1/2$:

\begin{itemize}
\item{} For {\bfseries\boldmath Part 1} ($S = \Lambda$, $|S| = V$): $\frac{\operatorname{ive}(0, \beta/2)}{\operatorname{ive}(0, \beta)} = \sqrt{2}(1 + O(\beta^{-1}))$, and by Lemma 3.0 the $n \ne 0$ sum adds at most a constant factor $C_W(e\ell)$ (since $(|n|-1)^2 V / (2(\beta/2 + |n|)) = (e\ell)^2 (|n|-1)^2 / (1 + 2|n|/\beta)$), giving $\mathbb{E}_W[e^{\frac{\beta}{2} \sum_p (1 - \cos\theta_p)}] \le C_W 2^{V/2}(1 + O(\beta^{-1}))^V \le e^{0.4 V}$. Markov at $K(e\ell)^2$ yields $\mu_W(\neg G_{\text{act}}) \le e^{-\frac{1}{2} K V + 0.4 V} \le e^{-(c_f(H_T) + 1) V}$.
\item{} For {\bfseries\boldmath Part 2} ($S = S_{\text{cut}}^+$, $|S| = O(L \ln L) \le V/2$): the denominator is $\ge \operatorname{ive}(0, \beta)^V$. For $n = 0$, the ratio is $(\operatorname{ive}(0, \beta/2)/\operatorname{ive}(0, \beta))^{|S|} \le 2^{|S|} = e^{(\ln 2) |S|}$; for $|n| \ge 1$, since $V - |S| \ge V/2$, Lemma 3.0 bounds the $n \ne 0$ sum by $C_W'(e\ell)$ times the $n = 0$ term. Thus

\begin{equation*}
\adjustbox{max width=\linewidth}{$\displaystyle \mathbb{E}_W\!\left[\exp\!\Big(\frac{\beta}{2} \sum_{p \in S_{\text{cut}}^+} (1 - \cos\theta_p)\Big)\right] \;\le\; C_W'(e\ell)\, e^{(\ln 2)|S_{\text{cut}}^+|}.$}
\end{equation*}

On $(G \cap G_{\text{act}}) \setminus E_{\text{strip}}$, Corollary 2b(4) gives $\frac{1}{2} L_5 C_{\text{PV}}(\mu_0) \sum_{p \in S_{\text{cut}}^+} \theta_p^2 \ge \frac{1}{2} L^{-3/4}$ (for large $L$), and on $G$, $1 - \cos\theta_p \ge \frac{1}{3} \theta_p^2$ (since $|\theta_p| \le \theta_\varepsilon < \pi/3$). Using $L_5 \le \frac{6}{\mu_T} \ln L$, we obtain $\sum_{p \in S_{\text{cut}}^+} (1 - \cos\theta_p) \ge \frac{\mu_T}{18 C_{\text{PV}}(\mu_0) L^{3/4} \ln L}$. Applying Markov with $\frac{\beta}{2} = \frac{L^2}{2(e\ell)^2}$ yields

\begin{equation*}
\adjustbox{max width=\linewidth}{$\displaystyle \mu_W\big((G \cap G_{\text{act}}) \setminus E_{\text{strip}}\big) \;\le\; C_W'(e\ell) \exp\!\left(-\frac{\mu_T L^{5/4}}{36(e\ell)^2 C_{\text{PV}}(\mu_0) \ln L} + (\ln 2)|S_{\text{cut}}^+|\right) \;\le\; \exp\!\left(-2 c_* \frac{L^{5/4}}{\ln L}\right)$}
\end{equation*}

for all sufficiently large $L$, with $c_* := \frac{\mu_T}{144(e\ell)^2 C_{\text{PV}}(\mu_0)} > 0$. $\blacksquare$
\end{itemize}


\subsection*{5. Proof of Main Theorem (\texttt{T1-OS'}, Route (v)) --- 4-Piece Assembly}
First record the two elementary free-torus \& Lane 1 bounds on $W_{\text{ov}}[U]$:

\begin{enumerate}
\item{} {\bfseries\boldmath Free AP torus Riemann-sum bound:} $\big|\ln \det D_{\text{ov}}^{H_T, (m_f)}[U_{\text{AP}}] + c_f(H_T) V\big| \le C_{\text{AP}} L$.
\emph{(Proof: On $U_{\text{AP}}$, $k \in \text{BZ}_{\text{AP}}$ is the midpoint grid of spacing $2\pi/L$ with $|k| \ge \sqrt{2}\pi/L$. First, $0 \le \ln((1-m_f^2)|\lambda_0|^2 + m_f^2) - \ln|\lambda_0|^2 \le m_f^2/|\lambda_0|^2$ sums to $O(\ln L)$, while $V \ln(1 - m_f^2) = O(1)$. Second, $|\nabla \ln|\lambda_0(q)|^2| \le C/|q|$ bounds the midpoint-rule error on cells $|k| \ge 2\pi/L$ by $\sum_k \frac{C}{L|k|} = O(L)$, and the 4 origin-adjacent cells contribute $O(\ln L)$.)}
\item{} {\bfseries\boldmath Global upper bound \& Lane 1 lower bound:} Since $D_{\text{ov}}^{H_T, (m_f)}[U]$ is normal with spectrum on $|z - \frac{1+m_f}{2}| = \frac{1-m_f}{2}$, $0 < \det D_{\text{ov}}^{H_T, (m_f)}[U] \le 1$ everywhere on $U(1)^{2V}$ (a.e. under $\mu_W$), so $W_{\text{ov}}[U] \le \exp(c_f(H_T) V + C_{\text{AP}} L)$ everywhere. Meanwhile, {\bfseries\boldmath Clawd's Lane 1 Lemma Z} (\texttt{LANE1-Z-LEMMA.md}, Section 8 for $H_T$ at fixed physical mass $m_f = \mu\ell/L$, signed off in Letter 89 and refereed in Letter 295 with $L_0(H_T) = 32$) gives $\ln \mathbb{E}_W[\det D_{\text{ov}}^{H_T, (m_f)}] \ge -c_f(H_T) V - O(L (\ln L)^{3/2})$ for every fixed $e\ell > 0$, hence

\begin{equation*}
\adjustbox{max width=\linewidth}{$\displaystyle \mathbb{E}_W[W_{\text{ov}}] \;=\; \frac{\mathbb{E}_W[\det D_{\text{ov}}^{H_T, (m_f)}]}{\det D_{\text{ov}}^{H_T, (m_f)}[U_{\text{AP}}]} \;\ge\; \exp\!\Big(-C_Z L (\ln L)^{3/2}\Big).$}
\end{equation*}

\end{enumerate}

Now partition $U(1)^{2V}$ into four disjoint regions and bound $\eta_L := \frac{\mathbb{E}_W[|W_{\text{ov}} - W_{\text{aux}}|]}{\mathbb{E}_W[W_{\text{ov}}]}$:

\begin{itemize}
\item{} {\bfseries\boldmath Piece A (Good cut strip $G \cap G_{\text{act}} \cap E_{\text{strip}}$):} By Corollary 2b(4), $|\ln W_{\text{aux}} - \ln W_{\text{ov}}| \le L^{-3/4} + O(L^{-2}) \le \frac{3}{2} L^{-3/4}$, so $|W_{\text{ov}} - W_{\text{aux}}| \le 2 L^{-3/4} W_{\text{ov}}$. Therefore $W_{\text{ov}}$ {\bfseries\boldmath cancels identically}:

\begin{equation*}
\adjustbox{max width=\linewidth}{$\displaystyle \frac{\mathbb{E}_W\big[|W_{\text{ov}} - W_{\text{aux}}|\, \mathbf{1}_{G \cap G_{\text{act}} \cap E_{\text{strip}}}\big]}{\mathbb{E}_W[W_{\text{ov}}]} \;\le\; 2 L^{-3/4}\, \frac{\mathbb{E}_W[W_{\text{ov}} \mathbf{1}_{G \cap G_{\text{act}} \cap E_{\text{strip}}}]}{\mathbb{E}_W[W_{\text{ov}}]} \;\le\; 2 L^{-3/4}.$}
\end{equation*}

\item{} {\bfseries\boldmath Piece B (Rare cut-strip spike $(G \cap G_{\text{act}}) \setminus E_{\text{strip}}$):} By Corollary 2b(2, 3), $|W_{\text{ov}} - W_{\text{aux}}| \le 4 L^{C_0}$. By Lemma 3(2) and Lane 1's lower bound on $\mathbb{E}_W[W_{\text{ov}}]$:

\begin{equation*}
\adjustbox{max width=\linewidth}{$\displaystyle \frac{\mathbb{E}_W\big[|W_{\text{ov}} - W_{\text{aux}}|\, \mathbf{1}_{(G \cap G_{\text{act}}) \setminus E_{\text{strip}}}\big]}{\mathbb{E}_W[W_{\text{ov}}]} \;\le\; 4 L^{C_0} \exp\!\left(-2 c_* \frac{L^{5/4}}{\ln L} + C_Z L (\ln L)^{3/2}\right) \;\le\; \frac{1}{2} \exp\!\left(-c_*\, \frac{L^{5/4}}{\ln L}\right)$}
\end{equation*}

(for large $L$, since $L^{5/4}/\ln L \gg L (\ln L)^{3/2}$!).
\item{} {\bfseries\boldmath Piece C (Action-cap complement $G \setminus G_{\text{act}}$):} Here $W_{\text{ov}} \le \exp(c_f(H_T) V + C_{\text{AP}} L)$ and $W_{\text{aux}} \le L^{C_0}$ (Corollary 2b(2)). By Lemma 3(1), $\mu_W(\neg G_{\text{act}}) \le e^{-(c_f(H_T) + 1) V}$, so

\begin{equation*}
\adjustbox{max width=\linewidth}{$\displaystyle \frac{\mathbb{E}_W\big[|W_{\text{ov}} - W_{\text{aux}}|\, \mathbf{1}_{G \setminus G_{\text{act}}}\big]}{\mathbb{E}_W[W_{\text{ov}}]} \;\le\; 2 \exp\!\Big(-V + C_{\text{AP}} L + C_Z L (\ln L)^{3/2}\Big) \;=\; e^{-\Omega(L^2)}.$}
\end{equation*}

\item{} {\bfseries\boldmath Piece D (Inadmissible set $\neg G$):} Here $W_{\text{ov}} \le \exp(c_f(H_T) V + C_{\text{AP}} L)$ and $W_{\text{aux}} \le L^{C_0}$ (Corollary 2b(2) holds \emph{everywhere}, including on $(G_+^{\text{cap}} \cap G_-^{\text{cap}}) \setminus G$!). By Lane 1 eq. (1.1), $\mu_W(\neg G) \le C' V^2 \exp(-\frac{\varepsilon^2}{2(e\ell)^2} V)$ (using $1 - \cos\theta_\varepsilon = \varepsilon^2 / 2$). Since $e\ell < e\ell_0 = \varepsilon / \sqrt{2 c_f(H_T)}$, $\delta_0 := \frac{\varepsilon^2}{2(e\ell)^2} - c_f(H_T) > 0$, giving

\begin{equation*}
\adjustbox{max width=\linewidth}{$\displaystyle \frac{\mathbb{E}_W\big[|W_{\text{ov}} - W_{\text{aux}}|\, \mathbf{1}_{\neg G}\big]}{\mathbb{E}_W[W_{\text{ov}}]} \;\le\; C'' V^2 \exp\!\Big(-\delta_0 V + C_{\text{AP}} L + C_Z L (\ln L)^{3/2}\Big) \;=\; e^{-\Omega(L^2)}.$}
\end{equation*}

\end{itemize}

Summing Pieces A--D (with Pieces C + D $\le \frac{1}{2} \exp(-c_* L^{5/4}/\ln L)$ for large $L$) yields $\eta_L \le 2 L^{-3/4} + \exp(-c_* L^{5/4}/\ln L) < 1$ for all sufficiently large $L$. In particular $\mathbb{E}_W[W_{\text{aux}}] \ge (1 - \eta_L) \mathbb{E}_W[W_{\text{ov}}] > 0$, and the elementary triangle inequality gives

\begin{equation*}
\adjustbox{max width=\linewidth}{$\displaystyle d_{\text{TV}}\!\left(\mu_{\text{ov}}^{(L)},\, \mu_{\text{aux}}^{(L)}\right) = \frac{1}{2} \mathbb{E}_W\!\left[\left|\frac{W_{\text{ov}}}{\mathbb{E}_W W_{\text{ov}}} - \frac{W_{\text{aux}}}{\mathbb{E}_W W_{\text{aux}}}\right|\right] \le \frac{1}{2}\left(\frac{\mathbb{E}_W[|W_{\text{ov}} - W_{\text{aux}}|]}{\mathbb{E}_W W_{\text{ov}}} + \left|1 - \frac{\mathbb{E}_W W_{\text{aux}}}{\mathbb{E}_W W_{\text{ov}}}\right|\right) \le \eta_L = O\!\big(L^{-3/4}\big).$}
\end{equation*}

Finally, since $\mathbb{E}_{\text{aux}}^{(L)}[(\Theta F) F] \ge 0$ (Corollary 2b(1)) and $|(\Theta F)(U) F(U)| \le \|F\|_\infty^2$,

\begin{equation*}
\adjustbox{max width=\linewidth}{$\displaystyle \mathbb{E}_{\text{ov}}^{(L)}\!\big[(\Theta F) F\big] \;\ge\; \mathbb{E}_{\text{aux}}^{(L)}\!\big[(\Theta F) F\big] - 2\, d_{\text{TV}}\!\left(\mu_{\text{ov}}^{(L)},\, \mu_{\text{aux}}^{(L)}\right) \|F\|_\infty^2 \;\ge\; -2\, d_{\text{TV}}\!\left(\mu_{\text{ov}}^{(L)},\, \mu_{\text{aux}}^{(L)}\right) \|F\|_\infty^2 \;=\; -O\!\big(L^{-3/4}\big) \|F\|_\infty^2. \qquad \blacksquare$}
\end{equation*}



\subsection*{6. Dependency Audit \& Audit Record (Letters 294 \& 295 Folded In)}
\begin{itemize}
\item{} {\bfseries\boldmath Reused from Clawd's \texttt{LANE1-Z-LEMMA.md} (refereed in Letters 290/89/295; dependency (13) CLOSED):}
\begin{itemize}
\item{} Step 1b eq. (1.1): $\mu_W(\neg G) \le C' V^2 \exp(-\frac{\varepsilon^2}{2(e\ell)^2} V)$.
\item{} Section 8 (\texttt{H\_T} Lemma Z, holding for every fixed $e\ell > 0$, with $L_0(H_T) = 32$ under $\eta_T = 0.0801$): $\ln \mathbb{E}_W[\det D_{\text{ov}}^{H_T, (m_f)}] \ge -c_f(H_T) V - O(L (\ln L)^{3/2})$.
\end{itemize}

\item{} {\bfseries\boldmath Proved 100\% inline (zero external citations to Neuberger, Adams, Osterwalder--Seiler, Simon, Borell--TIS, Caffarelli, or Pr\'ekopa):}
\begin{itemize}
\item{} Lemma 1 (strict positivity + exterior-power diamagnetic bound + exact link RP of $W_{\text{aux}}$).
\item{} Lemma 2.0 (5-line local 2D spectral gap $\mu_0^2(m, \varepsilon) = m^2 - \varepsilon - \frac{\varepsilon^2}{8} - \frac{\varepsilon^2}{2(1 - m - \varepsilon/4)}$, giving $\mu_0 = 0.24753$ at $m = 0.7, \varepsilon = 0.267$ and remaining strictly positive up to $\varepsilon_{\text{2D}}^*(0.7) = 0.291819$, holding locally on $B_r(x)$ using only plaquettes in $B_{r+1}(x) \subset B_{r+2}(x)$).
\item{} Lemma 2a \& Corollary 2b (exact 9-block Combes--Thomas with $C_H = 16(e^{1/2} + e)$ + 3-line Dirichlet Schur identity for all $x \in \Lambda$ on $G$ + corrected $\mathcal{C} M^T \mathcal{C}^{-1} = \gamma_5 M[U^*] \gamma_5$ site-trace vanishing + exact Hilbert-space Term 2 bound with explicit $C_3(\mu_0)$ and explicit $C_{\text{geom}}(\gamma_0) = \frac{4q(1+q)}{(1-q)^4}$ + ceiling-safe $C_0$ + global ceiling $W_{\text{aux}} \le L^{C_0}$ + strip bound on $\Delta_{\text{cut}}$).
\item{} Lemma 3.0 \& Lemma 3 (1-line Vandermonde proof of Tur\'an's inequality + exact plaquette-dependent 2D character expansion with explicit $c_* = \frac{\mu_T}{144(e\ell)^2 C_{\text{PV}}(\mu_0)}$).
\item{} Section 5 (4-line AP midpoint Riemann-sum bound + 4-piece assembly).
\end{itemize}

\item{} {\bfseries\boldmath Clawd's Letter 294 audit \& folded wounds (\texttt{2026-11-26-route5-audit-clawd/\allowbreak{}}):}
\begin{itemize}
\item{} {\bfseries\boldmath Ruling:} HOLDS conditionally on Lane 1 Section 8. Independent from-scratch code (\texttt{lat.py}, \texttt{audit\_checks.py}, \texttt{p2\_more.py}, \texttt{p4\_gap.py}) confirmed P0 (\texttt{1e-13}), P1 (diamagnetic ceiling across 160 Haar/soft + 252 flat non-AP fields), P2/P3 (7/7 $\beta = 0$ link-RP setups PSD for $\det D_{\text{DW}}$, while $\det D_{\text{ov}}$ negative control fails all 7/7), and P4 (gap floor \texttt{0.443} / \texttt{0.955} vs. \texttt{0.0613}). The 15-minute cold refuter found zero internal kills and further stress-tested Lemma 1(2) with L-BFGS ascent of $\ln\det D_{\text{DW}}$ on $L = 4$ and $L = 6$ (every ascent converged to $U_{\text{AP}}$ from below, never above).
\item{} {\bfseries\boldmath Wounds folded in rev. 2:}
\begin{enumerate}
\item{} \texttt{(1')} Fixed line 73 typo: $D_W(M) = (M + 2) I - (X_0 + X_1)$.
\item{} \texttt{(6')} Restated Lemma 2a(1) for {\bfseries\boldmath every $x \in \Lambda$} on $G$ (with explicit constant $896 r_0 / c_{\text{PV}}^2$), and derived the full-torus Lemma 2a(2) directly from Lemma 2a(1) + the Dirichlet-ball bound.
\item{} \texttt{(7)} Replaced the Term 2 estimate in Lemma 2a(2) with the exact Hilbert-space operator bound $\operatorname{tr}_{2\times 2}(\Psi^\dagger R_1 \Psi) \le \frac{2}{c_{\text{PV}}^2 + t}\sum_v \|\Psi_v\|_{\text{op}}^2$, preserving exact $e^{-2\gamma_0 d(x, l)}$ decay with zero $d^5/\gamma_0^2$ blow-up and explicit closed-form $C_3(\mu_0)$, $C_{\text{geom}}(\gamma_0)$.
\item{} \texttt{(10)} Dropped ``quantitative'' from title/theorem/prose and stated plainly that the bound is asymptotic in $L$ (crossover $L \sim 10^{138}\text{--}10^{141}$), alongside the exact scope and comparison to OpenAI \S{}125.
\item{} \texttt{(12)} Made $C_H, C_3(\mu_0), C_{\text{geom}}(\gamma_0), C_{\text{PV}}(\mu_0)$, and $c_*$ completely explicit (plus ceiling-safe $C_0 = 6 C_{\text{PV}} K'(e\ell)^2 / \mu_T$).
\end{enumerate}

\end{itemize}

\item{} {\bfseries\boldmath Clawd's Letter 295 Section 8 refuter round \& firsthand Neuberger/Adams verification (folded in rev. 3):}
\begin{itemize}
\item{} {\bfseries\boldmath Lane 1 Section 8 ruling:} Rate $\ln \mathbb{E}_W[\det D_{\text{ov}}^{H_T, (m_f)}] \ge -c_f(H_T) V - O(L(\ln L)^{3/2})$ HOLDS (all numerical checks, $c_f(H_T) = 0.1589390$, $\mu_T(m)$, L-BFGS saturation on $\|(2+D_W)^{-1}\| \le 1/(2-m)$, $H_{T1}$ gap, $h_T$, $n_0/T_0$, and $U_{\text{AP}}$ density survived); $L_0(H_T)$ raised to $32$ under $\eta_T = 0.0801$.
\item{} {\bfseries\boldmath Firsthand verification of Neuberger \& Adams (Row 11 CLOSED):} Direct reading of the arXiv TeX sources of H. Neuberger, \emph{Bounds on the Wilson Dirac Operator}, Phys. Rev. D {\bfseries\boldmath 61}, 085015 (\texttt{hep-lat/\allowbreak{}9911004}, \texttt{v2.tex} lines 506--614) and D. H. Adams, \emph{General bounds on the Wilson-Dirac operator}, Phys. Rev. D {\bfseries\boldmath 68}, 065009 (\texttt{hep-lat/\allowbreak{}9907005}, \texttt{newbwd.tex} eq. (2.1)) confirms that Neuberger proved $|H_W(-m)| \ge \sqrt{1 - (2+\sqrt{2})\sum_{\mu>\nu}\varepsilon_{\mu\nu}} - |1 - m|$ directly in arbitrary dimension $d$ ($r = 1$, $U(1)/SU(n)$), with $H_W(-1)^2 = 1 + Q + X + Y + Z$ verified numerically to \texttt{2.89e-15} at $d = 2$. Under Neuberger's bound at $m = 0.7$, $\varepsilon_{\text{Neu}}^* = 0.26653283$, giving open $e\ell < 0.4727$ ($H_T$) and $e\ell < 0.3367$ ($H_W$) --- confirming Letter 295's kill of Adams' bound at $\varepsilon = 0.267$.
\item{} {\bfseries\boldmath Inline Lemma 2.0 comparison:} Because Route (v)'s inline 5-line 2D Lemma 2.0 keeps $(1 - m)(\nabla_0^\dagger \nabla_0 + \nabla_1^\dagger \nabla_1)$ directly at $m < 1$, it strictly beats Neuberger at $m = 0.7$: at fixed $\varepsilon = 0.267$ it gives the uniform gap $\mu_0 = 0.24753 > 0$ (covering all $e\ell < 0.4727$ and up to $0.4735$ without gap degeneration), and its zero $\varepsilon_{\text{2D}}^* = 0.29181907$ extends the open 2D ceiling to $e\ell < 0.5175$ ($H_T$) / $e\ell < 0.3687$ ($H_W$) with zero external citation (all checks verified in \texttt{verify\_lane2\_route5\_bundle.py}).
\item{} {\bfseries\boldmath Pre-Release-11 dependencies:} None remaining; ready for Release 11 (\texttt{T1-OS'}) packaging.
\end{itemize}

\end{itemize}


\section*{Part II. LANE 1: the finite-beta lower bound on Z at fixed physical fermion mass (Clawd; started after clock reading 12:09, 2026-11-26)}
Status: SIGNED OFF by Wren (letter 89, 13:01): sections 0 to 7 HOLD as written. Section 8 held at m = 0.7 and needed
five small fixes for all 0 $<$ m $<$= 1; they are now in, together with her sharpening of the statement (the
V ln(1 - m\_f\textasciicircum{}2) term is gone). Before that, one refuter round on sections 0 to 7 returned HOLDS, with four minor
wounds, all repaired. A second refuter round, on Section 8 (23:56, 11-26), left the rate standing and killed one
constant: eps = 0.267 sat past the zero of Adams' bound, so the ceilings are now the open bounds e l $<$ 0.4727 (H\_T)
and e l $<$ 0.3367 (H\_W), and L\_0(H\_T) is 32. Section 9 has all three records. The numeric controls (sealed plan
\texttt{verification/\allowbreak{}lane1/\allowbreak{}lane1-plan.md}) test the inequalities and the normalisations on exact Wilson
samples, not the theorems cited.
Split and lead: Wren letter 87 part 3.2(b), my letter 289, Wren letter 88 (fixed PHYSICAL mass mu $>$ 0, so m\_f = mu a).


\subsection*{0. Setup and notation}
\begin{itemize}
\item{} Lattice: the L x L torus, V = L\textasciicircum{}2 sites, L $>$= 3. Links U\_l = e\textasciicircum{}\{i A\_l\}, plaquette angles theta\_p in (-pi, pi]
(principal value of U\_p).
\item{} Wilson measure: mu\_W(dU) proportional to prod\_p e\textasciicircum{}\{beta cos theta\_p\} prod\_l dU\_l, with beta = V/(e l)\textasciicircum{}2 (fixed
physical size l = L a, fixed e).
\item{} Kernel: H = H\_W[U] = gamma5 (D\_W[U] - m), with r = 1, 0 $<$ m $<$= 1 (m = 0.7 in every number below), and the
conventions of T-b/T-e (\texttt{te\_finite\_beta.py}). Fermion boundary conditions: any fixed twist (antiperiodic in T-c/T-e);
the result does not depend on it (Step 4).
\item{} Normalised massive overlap, Wren's normalisation (letter 88): D\_0[U] = (1/2)(1 + gamma5 sign H),
D[U] = D\_\{m\_f\}[U] = (1 - m\_f) D\_0[U] + m\_f.
\begin{itemize}
\item{} W := gamma5 sign H is unitary, so D\_0 = (1 + W)/2 and D are normal, with eigenvalues (1 - m\_f)(1 + w)/2 + m\_f,
$|$w$|$ = 1.
\item{} So every singular value of D lies in [m\_f, 1], and sigma\textasciicircum{}2 = (1 - m\_f\textasciicircum{}2)$|$lambda\_0$|$\textasciicircum{}2 + m\_f\textasciicircum{}2 eigenvalue by
eigenvalue, where lambda\_0 = (1 + w)/2.
\item{} gamma5 D gamma5 = D\textasciicircum{}dag, so det D is real. Its real eigenvalues are 1 and m\_f, so det D[U] $>$ 0 at m\_f $>$ 0 for
every U where H\_W[U] is invertible (sign H\_W needs that). The set \{det H\_W[U] = 0\} is the zero set of a
real-analytic function that is not identically zero. So it is mu\_W-null, and it does not meet E, where the gap is
$>$= m/2 (Wren, letter 89).
\item{} Also $|$lambda\_0$|$ $<$= 1, so sigma\textasciicircum{}2 - $|$lambda\_0$|$\textasciicircum{}2 = m\_f\textasciicircum{}2 (1 - $|$lambda\_0$|$\textasciicircum{}2) $>$= 0. Each singular value of D\_\{m\_f\}
is at least the matching $|$lambda\_0$|$ (Wren, letter 89; used in section 5).
\end{itemize}

\item{} Fixed physical mass: m\_f = mu a = mu l / L, with mu l $>$ 0 fixed, so m\_f -$>$ 0 and V m\_f\textasciicircum{}2 = (mu l)\textasciicircum{}2 is O(1).
\item{} Free constant, H\_W kernel: c\_f(m) := -(1/(4 pi\textasciicircum{}2)) int\_\{T\textasciicircum{}2\} d\textasciicircum{}2k ln[(1 + b/w)/2], with
b = sum\_mu (1 - cos k\_mu) - m, s\textasciicircum{}2 = sum\_mu sin\textasciicircum{}2 k\_mu, w = sqrt(s\textasciicircum{}2 + b\textasciicircum{}2).
c\_f(0.7) = 2 ln 2 - f\_inf = 0.3131624 (T-c, f\_inf = 1.0731320).
\item{} Z := E\_W[det D\_\{m\_f\}[U]].
\end{itemize}


\subsection*{1. Statement}
{\bfseries\boldmath Lemma Z.} Fix 0 $<$ m $<$= 1, e l $>$ 0, mu l $>$ 0 and eps0 in (0, pi/2). Then there is an explicit L\_0, and for every
L $>$= L\_0,

\begin{Verbatim}[breaklines=true,breakanywhere=true,fontsize=\small]
ln Z >= -c_f(m) V - Lambda(L) + ln P_E,
\end{Verbatim}

where P\_E $>$= (1 - eta\_C)(tanh(a\_Q/2))(2/3 - eta\_inf) $>$ 0 (Step 1) and Lambda(L) = O(L (ln L)\textasciicircum{}\{3/2\}) (Step 3, (3.6)).
So

\begin{Verbatim}[breaklines=true,breakanywhere=true,fontsize=\small]
(1/V) ln Z >= -c_f(m) - O(L^{-1} (ln L)^{3/2}).
\end{Verbatim}

The first version carried an extra + V ln(1 - m\_f\textasciicircum{}2), which is $>$= -(mu l)\textasciicircum{}2/(1 - m\_f\textasciicircum{}2). Wren's sharpening
sigma $>$= $|$lambda\_0$|$ (section 0) removes it.

The lemma holds at EVERY fixed e l. The small box enters only in its use (Section 6).


\subsection*{2. Step 1: the event E, and its probability JOINTLY}
E is defined on the plaquette angles only (Step 4 shows no holonomy window is needed at m\_f $>$ 0):

\begin{Verbatim}[breaklines=true,breakanywhere=true,fontsize=\small]
E = C cap {Q = 0} cap {||a||_2^2 <= A*} cap {||a||_inf <= eta},
\end{Verbatim}

\begin{itemize}
\item{} C = \{max\_p $|$theta\_p$|$ $<$= eps0\}, the plaquette cap.
\item{} Q = (1/2pi) sum\_p theta\_p, an integer.
\item{} On Q = 0, a = a(theta) is the Coulomb potential: the unique real 1-form with da = theta (as real numbers),
d*a = 0 and zero mean. a = d* G theta, with G = (d d*)\textasciicircum{}+ on 2-forms. In Fourier, $|$$|$a$|$$|$\_2\textasciicircum{}2 = $<$theta, G theta$>$ =
sum\_\{k != 0\} $|$theta\textasciicircum{}\_k$|$\textasciicircum{}2 / khat\textasciicircum{}2, with khat\textasciicircum{}2 = sum\_mu 4 sin\textasciicircum{}2(k\_mu/2) and theta\textasciicircum{}\_k = V\textasciicircum{}\{-1/2\} sum\_p e\textasciicircum{}\{-ikp\} theta\_p.
\item{} A* = 3 (e l)\textasciicircum{}2 G(L)/cos eps0 and eta = m/4, where G(L) := (1/V) sum\_\{k != 0\} 1/khat\textasciicircum{}2 (an explicit finite sum).
\end{itemize}

{\bfseries\boldmath 1a. Representation.} The map U -$>$ (U\_p) is a homomorphism onto \{prod\_p U\_p = 1\}, with kernel the flat fields.
Modulo gauge, the flat fields are a 2-torus of holonomies.

\begin{itemize}
\item{} So under Haar, with any measurable section, the configuration is (theta, phi), with phi uniform on T\textasciicircum{}2 and
independent of theta. The Wilson weight depends on theta alone, so phi stays Haar and independent under mu\_W.
\item{} On Q = 0 the section is A = a(theta) + phi/L (constant phi\_mu/L on mu-links). e\textasciicircum{}\{ia\} has plaquette angles exactly
theta, and U e\textasciicircum{}\{-ia\} is flat, hence gauge-equivalent to the constant holonomy field.
\item{} Law of theta under mu\_W: iid von Mises(beta), conditioned on sum\_p theta\_p in 2 pi Z.
\begin{itemize}
\item{} Equivalently, theta\_1..theta\_\{V-1\} have density proportional to prod\_\{p$<$V\} e\textasciicircum{}\{beta cos theta\_p\} e\textasciicircum{}\{beta cos S'\},
where S' = sum\_\{p$<$V\} theta\_p. This is because U\_V = (prod\_\{p$<$V\} U\_p)\textasciicircum{}\{-1\}, so cos theta\_V = cos S'.
\item{} This law is exchangeable in all V plaquettes.
\end{itemize}

\end{itemize}

{\bfseries\boldmath 1b. The cap.} For any event F in theta\_1..theta\_\{V-1\},

\begin{Verbatim}[breaklines=true,breakanywhere=true,fontsize=\small]
mu_W(F) = E_0[1_F e^{beta(cos S' - 1)}] / E_0[e^{beta(cos S' - 1)}] <= P_0(F) / (e^{-beta} I_0(beta)),
\end{Verbatim}

where E\_0 is iid vM(beta) on V - 1 plaquettes. The step uses e\textasciicircum{}\{beta cos x\} = sum\_n I\_n(beta) e\textasciicircum{}\{inx\} and
E\_0 e\textasciicircum{}\{inS'\} = (I\_n/I\_0)\textasciicircum{}\{V-1\} $>$= 0, so E\_0 e\textasciicircum{}\{beta cos S'\} $>$= I\_0.

\begin{itemize}
\item{} Also e\textasciicircum{}\{-beta\} I\_0(beta) = (1/pi) int\_0\textasciicircum{}pi e\textasciicircum{}\{-beta(1 - cos t)\} dt $>$= (1/pi) int\_0\textasciicircum{}pi e\textasciicircum{}\{-beta t\textasciicircum{}2/2\} dt
= erf(pi sqrt(beta/2))/sqrt(2 pi beta).
\item{} And P\_vM($|$theta$|$ $>$ eps0) $<$= (1 - eps0/pi) e\textasciicircum{}\{-beta(1 - cos eps0)\} / (e\textasciicircum{}\{-beta\} I\_0(beta)).
\item{} By exchangeability the bound for theta\_1 holds for every p, including p = V. So
\begin{Verbatim}[breaklines=true,breakanywhere=true,fontsize=\small]
mu_W(not C) <= eta_C := V (1 - eps0/pi) 2 pi beta e^{-beta(1 - cos eps0)} / erf(pi sqrt(beta/2))^2.     (1.1)
\end{Verbatim}

\end{itemize}

{\bfseries\boldmath 1c. Q = 0, uniformly in L.} Under mu\_W(. $|$ C), theta is iid capped vM conditioned on S = sum theta in 2 pi Z, so
mu\_W(Q = q $|$ C) = f(2 pi q) / sum\_n f(2 pi n), with f the density of S under iid capped vM.

\begin{itemize}
\item{} Capped vM is strongly log-concave: e\textasciicircum{}\{beta cos t\} 1\_\{$|$t$|$ $<$= eps0\} = g(t) e\textasciicircum{}\{-kappa t\textasciicircum{}2/2\}, with
kappa = beta cos eps0 and g log-concave and even. This holds because (beta cos t + kappa t\textasciicircum{}2/2)'' =
beta(cos eps0 - cos t) $<$= 0 on $|$t$|$ $<$= eps0.
\item{} Write theta = (x/V) 1 + theta\_perp, with theta\_perp in H = \{sum = 0\}. Then
\begin{Verbatim}[breaklines=true,breakanywhere=true,fontsize=\small]
f(x) proportional to e^{-kappa x^2/(2V)} J(x),   J(x) = int_H prod_p g(x/V + theta_perp,p) e^{-kappa |theta_perp|^2/2} dtheta_perp.
\end{Verbatim}

J is log-concave (Prekopa) and even, so J(x) $<$= J(0), and f(2 pi q)/f(0) $<$= e\textasciicircum{}\{-a\_Q q\textasciicircum{}2\}, with
a\_Q := kappa (2pi)\textasciicircum{}2/(2V) = 2 pi\textasciicircum{}2 cos eps0/(e l)\textasciicircum{}2.
\item{} Hence, with no dependence on L,
\begin{Verbatim}[breaklines=true,breakanywhere=true,fontsize=\small]
mu_W(Q = 0 | C) >= 1 / sum_n e^{-a_Q n^2} >= tanh(a_Q/2).                                               (1.2)
\end{Verbatim}

\end{itemize}

{\bfseries\boldmath 1d. The Coulomb norms under nu := mu\_W(. $|$ C, Q = 0).}

\begin{itemize}
\item{} nu is the slice measure on H $\cap$ [-eps0, eps0]\textasciicircum{}V, with density prod\_p e\textasciicircum{}\{beta cos theta\_p\}. Its potential has Hessian
$>$= kappa on a convex domain, and it is even.
\item{} Caffarelli's contraction theorem, used only with a FINITE convex perturbation, so no convex-support version is
needed.
\begin{itemize}
\item{} Let psi(t) = -beta cos t on $|$t$|$ $<$= eps0. Outside, take its C\textasciicircum{}1 quadratic continuation,
psi(t) = -beta cos eps0 + beta sin eps0 ($|$t$|$ - eps0) + (kappa/2)($|$t$|$ - eps0)\textasciicircum{}2, so psi'' $>$= kappa everywhere.
\item{} Let nu\_n be the measure on H with density proportional to
exp(-sum\_p psi(theta\_p) - n sum\_p ($|$theta\_p$|$ - eps0)\_+\textasciicircum{}3).
The cube makes the penalty C\textasciicircum{}2 and convex (Wren's polish, letter 89). psi is C\textasciicircum{}2 already, since its second
derivative is beta cos eps0 = kappa on both sides of eps0.
Its potential minus kappa $|$theta$|$\textasciicircum{}2/2 is a finite convex function on H.
\item{} So Caffarelli gives nu\_n = T\_n \# N(0, kappa\textasciicircum{}\{-1\} I\_H) with T\_n 1-Lipschitz. Hence Cov\_\{nu\_n\} $<$= kappa\textasciicircum{}\{-1\} I\_H
(Gaussian Poincare composed with T\_n), and nu\_n(f - E\_\{nu\_n\} f $>$= t) $<$= e\textasciicircum{}\{-kappa t\textasciicircum{}2/(2 Lip(f)\textasciicircum{}2)\}.
\item{} As n -$>$ infinity, nu\_n -$>$ nu weakly, by dominated convergence under the envelope e\textasciicircum{}\{-kappa $|$theta$|$\textasciicircum{}2/2\}.
The tails are uniformly Gaussian, so second moments converge.
\item{} So Cov\_nu $<$= kappa\textasciicircum{}\{-1\} I\_H. For the open sets \{$|$a\_l$|$ $>$ eta\}, nu $<$= liminf nu\_n (portmanteau). The two
consequences below hold for nu.
\end{itemize}

\item{} (l2) E\_nu theta = 0, so E\_nu $|$$|$a$|$$|$\textasciicircum{}2 = tr(G Cov\_nu) $<$= tr G / kappa = V G(L)/(beta cos eps0) = (e l)\textasciicircum{}2 G(L)/cos eps0.
By Markov, nu($|$$|$a$|$$|$\_2\textasciicircum{}2 $>$ A*) $<$= 1/3.
\item{} (sup) a\_l = $<$G d e\_l, theta$>$ is linear, with Lip\textasciicircum{}2 = $<$d e\_l, G\textasciicircum{}2 d e\_l$>$ = (1/V) sum\_\{k != 0\} 4 sin\textasciicircum{}2(k\_nu/2)/khat\textasciicircum{}4
= G(L)/2 by the square symmetry. E\_nu a\_l = 0. Taking a union over the 2V links,
\begin{Verbatim}[breaklines=true,breakanywhere=true,fontsize=\small]
nu(||a||_inf > eta) <= eta_inf := 4V exp(-kappa eta^2/G(L)) = 4V exp(-V cos eps0 eta^2/((e l)^2 G(L))).   (1.3)
\end{Verbatim}

\end{itemize}

{\bfseries\boldmath 1e. Jointly.} mu\_W(E) = mu\_W(C) mu\_W(Q = 0 $|$ C) nu(l2 and sup), so

\begin{Verbatim}[breaklines=true,breakanywhere=true,fontsize=\small]
P_E := mu_W(E) >= (1 - eta_C) tanh(a_Q/2) (2/3 - eta_inf).                                                (1.4)
\end{Verbatim}

The union bound is used only inside nu, where both pieces are small or Markov. Conditioning on C before Q is what
makes the Q bound L-uniform. A direct change of measure loses a factor sqrt(2 pi beta) \textasciitilde{} L, so it cannot bound
P(Q != 0) (checked by hand: P\_0($|$S'$|$ $>$= pi) does not decay in L).

{\bfseries\boldmath Explicit bound on G(L).} For k in (-pi, pi]\textasciicircum{}2, khat\textasciicircum{}2 $>$= 4$|$k$|$\textasciicircum{}2/pi\textasciicircum{}2. The shell $|$n$|$\_inf = j has 8j points with
$|$n$|$\textasciicircum{}2 $>$= j\textasciicircum{}2, so G(L) $<$= (1/16) sum\_\{0 $<$ $|$n$|$\_inf $<$= L/2\} $|$n$|$\textasciicircum{}\{-2\} $<$= (1 + ln(L/2))/2. Exact values (out-consts.log):
G(L) = ln L/(2 pi) + 0.04876 + o(1), for example G(16) = 0.4899 and G(64) = 0.7107.

Numbers (m = 0.7, eps0 = 0.5, eta = 0.175; out-consts.log, with eta\_inf evaluated using the proved bound on G):

\begin{center}
\begin{adjustbox}{max width=\linewidth}
\begin{tabular}{llll}
\hline
{}{\bfseries\boldmath e l} & {\bfseries\boldmath tanh(a\_Q/2)} & {\bfseries\boldmath eta\_C at L = 12 / 16 / 32} & {\bfseries\boldmath eta\_inf at L = 16 / 32 / 64 / 128} \\
\hline
{}0.34 & 1 - 2e-65 & 6e-61 / 6e-112 / 0 & 1.7e-14 / 6e-52 / 1e-181 / 0 \\
{}0.47 & 1 - 2e-34 & 1.1e-29 / 3.8e-56 / 9e-240 & 1.7e-6 / 8e-26 / 2e-93 / 0 \\
{}2 & 0.974 & vacuous / vacuous / 3.4e-8 & vacuous / vacuous / 0.073 / 1.9e-14 \\
{}4 & 0.494 & vacuous to L = 32 & vacuous to L = 128 \\
\hline
\end{tabular}
\end{adjustbox}
\end{center}

So at e l $<$= 0.47, P\_E $>$= 2/3 - 2e-6 for every L $>$= 16. Both tails decrease in L from L = 16 on.

\begin{itemize}
\item{} ln eta\_C = 2 ln V - V(1 - cos eps0)/(e l)\textasciicircum{}2 + O(1), where the 1/erf\textasciicircum{}2 factor decreases. This falls once
V $>$ 2 (e l)\textasciicircum{}2/(1 - cos eps0), which is V $>$ 3.7 at e l = 0.47.
\item{} ln eta\_inf = ln 4V - V cos eps0 eta\textasciicircum{}2/((e l)\textasciicircum{}2 G(L)), and V/G(L) grows faster than ln V. The exponent at L = 16 is
already 20. At e l = 2 the proved bound needs L $>$\textasciitilde{} 64, and at e l = 4,
L $>$\textasciitilde{} 200: L\_0 grows with e l, but P\_E stays Omega(1) at every fixed e l.
\end{itemize}


\subsection*{3. Step 2: perturbation, on E, uniformly in phi}
Let D1 = D[a + phi/L] and D0 = D[phi/L] (free, holonomy phi), with H1 and H0 their kernels.

\begin{itemize}
\item{} Free gap: H0\textasciicircum{}2 = s\textasciicircum{}2 + b\textasciicircum{}2 $>$= m\textasciicircum{}2 + 2 U(1 - m) $>$= m\textasciicircum{}2, for every real k (so every phi), where U = sum(1 - cos k\_mu).
This uses sin\textasciicircum{}2 = u(2 - u) with u = 1 - cos, and sum u\textasciicircum{}2 $<$= U\textasciicircum{}2.
\item{} Hopping difference: D\_W[A1] - D\_W[A0] = -sum\_\{x,mu\} [P\_-\textasciicircum{}mu Delta\_\{(x,mu)\} psi(x + mu) + P\_+\textasciicircum{}mu conj(Delta\_\{(x-mu,mu)\}) psi(x - mu)],
with Delta\_l = U1\_l - U0\_l, $|$Delta\_l$|$ = $|$e\textasciicircum{}\{i a\_l\} - 1$|$ $<$= $|$a\_l$|$, and P\_-\textasciicircum{}mu + P\_+\textasciicircum{}mu = 1 orthogonal projectors.
\begin{itemize}
\item{} Operator norm: per direction, (P\_- X + P\_+ Y)\textasciicircum{}dag (P\_- X + P\_+ Y) = P\_- X\textasciicircum{}dag X + P\_+ Y\textasciicircum{}dag Y, which is $<$=
max$|$Delta$|$\textasciicircum{}2. So $|$$|$H1 - H0$|$$|$\_op $<$= 2 $|$$|$a$|$$|$\_inf $<$= m/2 on E.
\item{} Hence the gap of H1 is $>$= m - m/2, and both kernels have spectra outside (-mu0, mu0), with mu0 = m/2.
\item{} HS norm: for L $>$= 3, each ordered neighbour pair carries one term, so $|$$|$H1 - H0$|$$|$\_HS\textasciicircum{}2 = 2 sum\_l $|$Delta\_l$|$\textasciicircum{}2 $<$=
2 $|$$|$a$|$$|$\_2\textasciicircum{}2.
\end{itemize}

\item{} Lipschitz in HS: for Hermitian A, B and Lipschitz f, $|$$|$f(A) - f(B)$|$$|$\_HS $<$= Lip(f) $|$$|$A - B$|$$|$\_HS. This is the
double sum sum\_ij $|$f(a\_i) - f(b\_j)$|$\textasciicircum{}2 $|$$<$u\_i, v\_j$>$$|$\textasciicircum{}2.
\begin{itemize}
\item{} Take f = sign off (-mu0, mu0), linear inside, so Lip = 1/mu0 = 2/m. Then
\begin{Verbatim}[breaklines=true,breakanywhere=true,fontsize=\small]
||D1 - D0||_HS = ((1 - m_f)/2) ||sign H1 - sign H0||_HS <= sqrt2 ||a||_2 / m =: h <= h_E := sqrt(6 (e l)^2 G(L)/cos eps0) / m.   (2.1)
\end{Verbatim}

\item{} h\_E\textasciicircum{}2 = O(ln L). At e l = 0.47, L = 1024: h\_E = 1.88.
\end{itemize}

\end{itemize}


\subsection*{4. Step 3: Mirsky with a mass floor (Wren's three groups, constants explicit)}
\begin{itemize}
\item{} Mirsky: sum\_i (sigma\_i(D1) - sigma\_i(D0))\textasciicircum{}2 $<$= $|$$|$D1 - D0$|$$|$\_F\textasciicircum{}2 $<$= h\textasciicircum{}2, with both sorted decreasingly, N = 2V.
\item{} Write Delta\_i = sigma\_i\textasciicircum{}1 - sigma\_i\textasciicircum{}0. Every sigma lies in [m\_f, 1]. Fix delta in (0, 1/4].
\item{} Moved, B = \{$|$Delta\_i$|$ $>$ delta\}: $|$B$|$ $<$= h\textasciicircum{}2/delta\textasciicircum{}2, and each costs at most ln(1/m\_f).
\item{} Small, S = \{i not in B, sigma\_i\textasciicircum{}0 $<$ 2 delta\}: each costs at most ln(1/m\_f), and $|$S$|$ $<$= n\_0(2 delta).
\item{} Good, the rest: sigma\textasciicircum{}1/sigma\textasciicircum{}0 $>$= 1/2, and ln(1 - x) $>$= -2 ln2 x on [0, 1/2]. So the good group costs at most
2 ln2 sum $|$Delta\_i$|$/sigma\_i\textasciicircum{}0 $<$= 2 ln2 h T\_0(2 delta)\textasciicircum{}\{1/2\}, with T\_0(t) := sum\_\{sigma\_i\textasciicircum{}0 $>$= t\} (sigma\_i\textasciicircum{}0)\textasciicircum{}\{-2\}.
\begin{Verbatim}[breaklines=true,breakanywhere=true,fontsize=\small]
ln det D1 - ln det D0 >= -[h^2/delta^2 + n_0(2 delta)] ln(1/m_f) - 2 ln2 h T_0(2 delta)^{1/2}.               (3.1)
\end{Verbatim}

\end{itemize}

Free counts, uniform in phi. Let C2 = 1/(1 - m\_f\textasciicircum{}2) and c1 = pi sqrt((2 + m\textasciicircum{}2) C2).

\begin{itemize}
\item{} If b $>$= 0, then $|$lambda\_0$|$\textasciicircum{}2 = (1 + b/w)/2 $>$= 1/2.
\item{} If b $<$ 0, then $|$k\_mu$|$ $<$ pi/2 (because m $<$= 1), and $|$lambda\_0$|$\textasciicircum{}2 = s\textasciicircum{}2/(2w(w - b)) $>$= s\textasciicircum{}2/(4(s\textasciicircum{}2 + m\textasciicircum{}2)) $>$=
$|$k$|$\textasciicircum{}2/(pi\textasciicircum{}2(2 + m\textasciicircum{}2)).
\item{} So for t $<$= 1/2 and m\_f\textasciicircum{}2 $<$ 1/2: sigma\textasciicircum{}0(k) $<$ t forces $|$k$|$ $<$ c1 t. Each k carries 2 singular values, and a disc of
radius r holds at most (rL/pi + 1)\textasciicircum{}2 grid momenta. Hence
\begin{Verbatim}[breaklines=true,breakanywhere=true,fontsize=\small]
n_0(2 delta) <= 16 (2 + m^2) C2 delta^2 V + 4.                                                                  (3.2)
\end{Verbatim}

\item{} Split T\_0(2 delta) at $|$k$|$ = 2 c1 delta (Wren's expository line, letter 89).
\begin{itemize}
\item{} Inner disc: at most the (3.2) count of singular values, each $>$= 2 delta. So it contributes at most
(16 (2 + m\textasciicircum{}2) C2 delta\textasciicircum{}2 V + 4)/(4 delta\textasciicircum{}2) = 4 C2 (2 + m\textasciicircum{}2) V + delta\textasciicircum{}\{-2\}. That is the second term of tau and
the + delta\textasciicircum{}\{-2\}.
\item{} Outer cells, in the annulus c1 delta $<$= $|$q$|$ $<$= sqrt2 pi (1 + 1/L): these give the 4 C2 and log terms, as follows.
\end{itemize}

\item{} With sigma\textasciicircum{}\{-2\} $<$= C2 [2 + pi\textasciicircum{}2 (2 + m\textasciicircum{}2)/$|$k$|$\textasciicircum{}2], the cell comparison ($|$q$|$ $<$= (3/2)$|$k$|$ on the cell of any grid k
with $|$k$|$ $>$= 2 sqrt2 pi/L) gives, for delta $>$= sqrt2 pi/(c1 L),
\begin{Verbatim}[breaklines=true,breakanywhere=true,fontsize=\small]
T_0(2 delta) <= V tau(delta) + delta^{-2},
tau(delta) = 4 C2 + 4 C2 (2 + m^2) + (9 pi/4) C2 (2 + m^2) ln(sqrt2 pi (1 + 1/L)/(c1 delta)).                   (3.3)
\end{Verbatim}

At m = 0.7 and small m\_f: tau \textasciitilde{} 14.0 + 17.6 ln(0.896/delta).
\item{} Choose delta\textasciicircum{}4 = h\_E\textasciicircum{}2/(16 (2 + m\textasciicircum{}2) C2 V), which balances the moved and small groups. Then
\begin{Verbatim}[breaklines=true,breakanywhere=true,fontsize=\small]
Lambda(L) := 2 sqrt(16 (2 + m^2) C2) h_E L ln(1/m_f) + 4 ln(1/m_f) + 2 ln2 h_E (V tau(delta) + delta^{-2})^{1/2}   (3.6)
         = O(L (ln L)^{3/2}),
\end{Verbatim}

and on E, for every phi, ln det D1 $>$= ln det D0 - Lambda(L).
\item{} Numbers at m = 0.7, e l = 0.47, mu l = 1:
\begin{itemize}
\item{} L = 256: Lambda/V = 0.54. L = 1024: delta = 0.017, Lambda/V = 0.18. L = 10\textasciicircum{}4: Lambda/V = 0.028
(\texttt{verification/\allowbreak{}lane1/\allowbreak{}lambda\_bound.py}, out-lambda.log).
\end{itemize}

\item{} The rate is honest but slow: Lambda/V is above c\_f = 0.313 until L \textasciitilde{} 500. The constants are crude (the small and moved
groups are charged the full ln(1/m\_f)). T-e's measured Delta is O(1) at L $<$= 24.
\end{itemize}


\subsection*{5. Step 4 and assembly: Jensen over the Haar holonomy (why no window is needed)}
\begin{itemize}
\item{} Free part, exact. For fixed n, k\_n = (2 pi n + chi + phi)/L, with chi the boundary twist, is uniform on one cell
when phi is uniform on T\textasciicircum{}2, and the L\textasciicircum{}2 cells tile T\textasciicircum{}2 once. With g(k) = ln$|$lambda\_0(k)$|$\textasciicircum{}2 in L\textasciicircum{}1:
\begin{Verbatim}[breaklines=true,breakanywhere=true,fontsize=\small]
E_phi ln det D_0[phi/L] = sum_n E_phi g(k_n) = (V/(4 pi^2)) int_{T^2} g = -c_f V   EXACTLY, for any twist chi.
\end{Verbatim}

The periodic-holonomy zero mode is an integrable log singularity and costs nothing in this average.
\item{} Mass. sigma $>$= $|$lambda\_0$|$ singular value by singular value (section 0), so ln det D\_\{m\_f\} $>$= ln$|$det D\_0$|$ for every
field. The first version used sigma\textasciicircum{}2 $>$= (1 - m\_f\textasciicircum{}2)$|$lambda\_0$|$\textasciicircum{}2 and paid V ln(1 - m\_f\textasciicircum{}2). Step 3's C2 could be 1
by the same inequality; it is left as is, which is valid and slightly loose.
\item{} Assembly. phi is independent of theta and Haar (1a). Exp is convex (Jensen). On E, (3.1) holds uniformly in phi.
Hence
\begin{Verbatim}[breaklines=true,breakanywhere=true,fontsize=\small]
Z >= E_W[det D; E] = int_E mu_W(dtheta) E_phi det D1 >= int_E mu_W(dtheta) exp(E_phi ln det D1)
  >= P_E exp(-c_f V - Lambda(L)).
\end{Verbatim}

That is Lemma Z.
\item{} L\_0 requires: L $>$= 3; m\_f = mu l/L $<$= 1/2; delta $<$= 1/4 and delta $>$= sqrt2 pi/(c1 L); eta\_C $<$ 1 and
eta\_inf $<$ 2/3.
\end{itemize}


\subsection*{6. Use: the small-box ceiling}
\begin{itemize}
\item{} Bad set: not G = \{some p with $|$1 - U\_p$|$ $>$= eps\}, that is $|$theta\_p$|$ $>$= theta\_eps = 2 arcsin(eps/2), so
1 - cos theta\_p $>$= eps\textasciicircum{}2/2.
\item{} By 1b, mu\_W(not G) $<$= V 2 pi beta e\textasciicircum{}\{-beta eps\textasciicircum{}2/2\}/erf(pi sqrt(beta/2))\textasciicircum{}2, and sup det D $<$= 1. So
\begin{Verbatim}[breaklines=true,breakanywhere=true,fontsize=\small]
sup_{not G} det x mu_W(not G) / Z <= exp(V [c_f - eps^2/(2 (e l)^2)] + Lambda(L) + O(ln V)),
\end{Verbatim}

which tends to 0 if e l $<$ eps/sqrt(2 c\_f).
\begin{itemize}
\item{} That is the ceiling of THIS method, which charges the bad set sup det $<$= 1. It is not a physical threshold.
\item{} With eps $<$ eps* = 0.266533 and c\_f = 0.31316, the H\_W ceiling is the open bound e l $<$ 0.3367 (eps*/sqrt(2 c\_f) =
0.336784), below Wren's 0.34. CORRECTED 23:57, 11-26: the earlier ``eps = 0.267, e l\_0 = 0.337'' rounded eps UP past
the zero of the bound it feeds (refuter round 2, section 9).
\item{} eps* comes from Neuberger's bound (hep-lat/9911004, PRD 61, 085015; v2.tex L592-618, read first-hand 11-27 by
Wren, letter 95, and by me, the arXiv TeX source v2.tex):
$|$H(m)$|$ $>$= sqrt(1 - (2 + sqrt2) sum\_\{mu $>$ nu\} eps\_\{mu nu\}) - $|$1 - m$|$ under $|$$|$1 - U\_\{mu nu\}(x)$|$$|$ $<$= eps\_\{mu nu\}, in
every dimension (``holds in both even and odd dimensions''). At d = 2 there is one plane, so c\_1 = 2 + sqrt2, positive
iff eps $<$ eps* = m(2 - m)/c\_1 (= 0.266533 at m = 0.7). Adams (hep-lat/9907005, L238-246) cites it and extends it to
the doubler regions. The earlier line called it ``Adams' generalisation of a 4d derivation''; it is Neuberger's own,
in general d.
\item{} So e l\_0 moves with eps. The refuter's uniform-flux probe (L = 20, $|$1 - U\_p$|$ up to 0.50) never brought
min$|$eig H\_W$|$ below 0.458. That is no worst-case check.
\end{itemize}

\item{} The 0.47 ceiling (H\_T kernel at a5 = 1, c\_f = 0.158939) is section 8.
\end{itemize}


\subsection*{7. Against Wren's lead (letter 87, 3.2(b)), and what is left}
\begin{itemize}
\item{} Same skeleton: an Omega(1) event, the HS perturbation, Hoffman-Wielandt/Mirsky with a mass floor, three groups.
Her remainder terms are recovered with explicit constants.
\item{} Differences, all mine, each held under one refuter round:
\begin{enumerate}
\item{} No holonomy window. Jensen over Haar phi makes the free term exact (Step 4).
\item{} The H\_W gap on E comes from a sup cap on the Coulomb field ($|$$|$a$|$$|$\_inf $<$= m/4) plus the free gap m, not from the
admissibility bound. So eps0 only has to be $<$ pi/2, for strong log-concavity.
\item{} P(Q = 0) is bounded uniformly in L by Prekopa on the capped law (1c). Her ``P(Q = 0) $>$= c(e l)'' was assumed.
\item{} E\_nu $|$$|$a$|$$|$\textasciicircum{}2 is bounded under the true (capped, Q = 0) law by Caffarelli/Brascamp-Lieb, not by Gaussian
analogy.
\end{enumerate}

\item{} Cited, not reproved: Prekopa (marginals of log-concave); Caffarelli's contraction for e\textasciicircum{}\{-Phi\} gamma with Phi
convex and finite (1d now needs only this case); Mirsky's singular-value inequality. All are standard. None has been
read first-hand in this lane.
\item{} Not covered: m\_f = 0 (the floor ln(1/m\_f) is used in two groups); P = det D\_N\textasciicircum{}(S) at finite S (route iii's
replacement); a better rate than L (ln L)\textasciicircum{}\{3/2\}.
\end{itemize}


\subsection*{8. The H\_T kernel at a5 = 1 (the 0.47 ceiling): Steps 2 and 3 with new constants (checked numerically, sealed; refuter round 2: rate holds, one constant corrected)}
Kernel: H\_T := gamma5 D\_W(-m) (2 + D\_W(-m))\textasciicircum{}\{-1\}.

\begin{itemize}
\item{} Hermitian: the two factors commute, and D\_W is gamma5-hermitian.
\item{} Free: $|$H\_T(k)$|$\textasciicircum{}2 = (b\textasciicircum{}2 + s\textasciicircum{}2)/((2 + b)\textasciicircum{}2 + s\textasciicircum{}2), and gamma5 H\_T(k) = b\_T + i s\_T.gamma, with
b\_T = (b(2 + b) + s\textasciicircum{}2)/n, s\_T = 2s/n, n = (2 + b)\textasciicircum{}2 + s\textasciicircum{}2.
\item{} c\_f(H\_T) = 0.158939: even midpoint grids N = 512 to 4096, O(N\textasciicircum{}-2) Richardson limit 0.15893896. This is Wren's
0.1589390 (letter 89), and T-d's 512\textasciicircum{}2 value was 0.158936.
\begin{itemize}
\item{} The first value, 0.158943, came from an odd n = 4001 midpoint grid. That grid puts a point at $|$k$|$ = 3e-17, where
ln$|$lambda0$|$\textasciicircum{}2 = -75, and the point alone adds 4.7e-6 (\texttt{verification/\allowbreak{}lane1/\allowbreak{}wren09/\allowbreak{}out-check.log}).
\end{itemize}

\item{} e l\_0 = eps*/sqrt(2 x 0.158939) = 0.472738, so the ceiling is the OPEN bound e l $<$ 0.4727 (rounded down). CORRECTED
23:57, 11-26 (refuter round 2, section 9): the earlier 0.4735 (and 0.4736 in letter 290) used eps = 0.267, which is past
eps* = m(2 - m)/(2 + sqrt2) = 0.266533, the zero of Adams' bound at m = 0.7 (at 0.267 the bound is -0.00267). Any
e l $<$ 0.4727 works: take eps in (e l sqrt(2 c\_f), eps*); the gap on G is then positive but small near the top.
\item{} Adams' eps $<$ eps* also gaps H\_T (Wren, letter 89).
\begin{itemize}
\item{} $|$$|$2 + D\_W(-m)$|$$|$\_op $<$= (4 - m) + 2 = 6 - m for every U.
\item{} So $|$H\_T$|$ = $|$H\_W (2 + D\_W(-m))\textasciicircum{}\{-1\}$|$ $>$= $|$H\_W$|$/(6 - m) $>$ 0 on the same admissible set $|$1 - U\_p$|$ $<$ eps.
\end{itemize}

\end{itemize}

(a) Free gap, every holonomy, 0 $<$ m $<$= 1.

\begin{itemize}
\item{} Put X = b\textasciicircum{}2 + s\textasciicircum{}2 $>$= m\textasciicircum{}2 + 2U(1 - m) (section 3), and (2 + b)\textasciicircum{}2 + s\textasciicircum{}2 = X + 4(1 - m + U).
\item{} So $|$H\_T$|$\textasciicircum{}2 $>$= 1/(1 + f(U)), with f(U) = 4(1 - m + U)/(m\textasciicircum{}2 + 2U(1 - m)) and U in [0, 4].
\item{} f' has the sign of m\textasciicircum{}2 - 2(1 - m)\textasciicircum{}2. So f is monotone on [0, 4], and f $<$= max(f(0), f(4)). That gives
mu\_T(m) := min( m/(2 - m), sqrt((m\textasciicircum{}2 + 8(1 - m))/(m\textasciicircum{}2 + 8(1 - m) + 4(5 - m))) ).
\begin{itemize}
\item{} The first version used f(4) only. That is wrong for m $<$ 2 - sqrt2 = 0.586, where f decreases: at m = 0.5,
min$|$H\_T$|$ = 1/3 against f(4)'s 0.437 (Wren, letter 89; reproduced in out-check.log, where min$|$H\_T$|$ = m/(2 - m)
exactly at m = 0.3 and 0.5).
\end{itemize}

\item{} At m = 0.7, f increases (0.49 $>$ 0.18): mu\_T\textasciicircum{}2 $>$= 2.89/20.09, so mu\_T $>$= 0.3792 (0.37928, rounded DOWN; the first
version's 0.3793 rounded this lower bound up). The numerical minimum is 0.3939.
\end{itemize}

(b) Perturbation.

\begin{itemize}
\item{} H\_T = gamma5 (1 - 2 (2 + D\_W)\textasciicircum{}\{-1\}), so H\_T1 - H\_T0 = 2 gamma5 (2 + D\_W1)\textasciicircum{}\{-1\} (D\_W1 - D\_W0) (2 + D\_W0)\textasciicircum{}\{-1\}, with
D\_W1, D\_W0 the Wilson operators D\_W(-m) at A1, A0 (not the overlaps D1, D0 of Steps 2 and 3).
\item{} 2 + D\_W(-m) = D\_W(2 - m), whose hermitian part is $>$= 2 - m (the Wilson term is $>$= 0 and the naive part is
antihermitian), so $|$$|$(2 + D\_W)\textasciicircum{}\{-1\}$|$$|$ $<$= 1/(2 - m) for every U.
\item{} Hence $|$$|$H\_T1 - H\_T0$|$$|$\_op $<$= (4/(2 - m)\textasciicircum{}2) $|$$|$a$|$$|$\_inf and $|$$|$H\_T1 - H\_T0$|$$|$\_HS $<$= (2 sqrt2/(2 - m)\textasciicircum{}2) $|$$|$a$|$$|$\_2.
\item{} Sup cap eta\_T := mu\_T (2 - m)\textasciicircum{}2/8 = 0.0801 (0.080123, rounded down) gives a gap $>$= mu\_T/2 for H\_T1.
\item{} Then h\_T = (1/2)(2/mu\_T)(2 sqrt2/(2 - m)\textasciicircum{}2) $|$$|$a$|$$|$\_2 $<$= 4.42 $|$$|$a$|$$|$\_2. The exact value is 4.4126; the first version
rounded this upper bound down to 4.41. Compare H\_W: sqrt2/m = 2.0203, so $<$= 2.03 $|$$|$a$|$$|$\_2.
\item{} (1.3) changes only through eta: eta\_inf,T = 4V exp(-V cos eps0 eta\_T\textasciicircum{}2/((e l)\textasciicircum{}2 G(L))).
\end{itemize}

(c) Free counts.

\begin{itemize}
\item{} b\_T $>$= 0 gives $|$lambda0$|$\textasciicircum{}2 $>$= 1/2.
\item{} b\_T $<$ 0 forces b $<$ 0, since 2 + b $>$ 0, so $|$k\_mu$|$ $<$ pi/2. It also forces s\textasciicircum{}2 $<$ -b(2 + b), so
n $<$ (2 + b)\textasciicircum{}2 - b(2 + b) = 2(2 + b) $<$ 4 for every 0 $<$ m $<$= 1 (Wren, letter 89; the first version had n $<$ 5 from
the m = 0.7 bound s\textasciicircum{}2 $<$= 0.91).
\begin{itemize}
\item{} Then $|$lambda0$|$\textasciicircum{}2 $>$= s\_T\textasciicircum{}2/(4 w\_T\textasciicircum{}2) $>$= s\_T\textasciicircum{}2/4 = s\textasciicircum{}2/n\textasciicircum{}2 $>$= s\textasciicircum{}2/16 $>$= $|$k$|$\textasciicircum{}2/(4 pi\textasciicircum{}2).
\item{} That uses w\_T = $|$H\_T$|$ $<$= 1, which holds for every k and m, because 1 - m + U $>$= 0. The first version said ``$<$ 1
since b $>$ -1'', which fails at m = 1, k = 0.
\item{} Numerically the minimum of $|$lambda0$|$\textasciicircum{}2/$|$k$|$\textasciicircum{}2 is 0.67, against the analytic 0.0253: lossy, but explicit.
\end{itemize}

\item{} General form, if $|$lambda0$|$\textasciicircum{}2 $>$= c$|$k$|$\textasciicircum{}2 where b $<$ 0: c1\textasciicircum{}2 = C2/c, n\_0(2 delta) $<$= 16 c1\textasciicircum{}2 delta\textasciicircum{}2 V/pi\textasciicircum{}2 + 4, and
tau(delta) = 4 C2 + 4 c1\textasciicircum{}2/pi\textasciicircum{}2 + (9/(4 pi)) c1\textasciicircum{}2 ln(sqrt2 pi (1 + 1/L)/(c1 delta)).
\begin{itemize}
\item{} H\_W has c = 1/(pi\textasciicircum{}2 (2 + m\textasciicircum{}2)), which recovers (3.2) and (3.3).
\item{} H\_T has c = 1/(4 pi\textasciicircum{}2), so 16 c1\textasciicircum{}2/pi\textasciicircum{}2 = 64 C2. The sealed check below used the weaker c = 4/(25 pi\textasciicircum{}2)
(100 C2), which also holds.
\end{itemize}

\end{itemize}

So Lemma Z holds for the H\_T overlap at every 0 $<$ m $<$= 1, with mu\_T(m) from (a), c\_f(H\_T) = 0.158939 at m = 0.7, and
the constants above. Lambda keeps the same form, with bigger constants. Audited by Wren (letter 89), with five fixes,
all now in.

\begin{itemize}
\item{} Where the H\_W numbers do not carry over (refuter round 2):
\begin{itemize}
\item{} P\_E: the H\_T cap eta\_T = 0.0801 replaces eta = 0.175 in Section 2. At e l = 0.47, eta\_inf,T is 14.7 at L = 16 (no
bound), 0.50 at L = 24 and 0.004 at L = 32. So section 2's ``P\_E $>$= 2/3 - 2e-6 for every L $>$= 16'' becomes
L $>$= 32 for H\_T.
\item{} Lambda(L)/V stays above c\_f(H\_T) until L is about 4500 (about 500 for H\_W). The order O(L (ln L)\textasciicircum{}\{3/2\}) is
unchanged; the bound only bites at large L.
\end{itemize}

\end{itemize}

Sealed check (\texttt{verification/\allowbreak{}lane1/\allowbreak{}ht/\allowbreak{}ht-plan.md}, sha f75d26d5e2187863, sealed at clock reading 12:31:52; run \texttt{ht\_check.py}, log
\texttt{verification/\allowbreak{}lane1/\allowbreak{}ht/\allowbreak{}out-ht-2026-11-26.log}):

\begin{itemize}
\item{} Every killer (Q1 to Q6) held in all six cells: e l in \{0.47, 2\}, L in \{8, 12, 16\}.
\item{} The one guess, on the ratio of the paired difference to the H\_W value, MISSED. That was a guess, not a killer.
\end{itemize}


\subsection*{9. Refuter record}
\begin{itemize}
\item{} Round 1 covered sections 0 to 7 and returned at clock reading 12:41. Verdict: HOLDS, with four minor wounds, none of
them load-bearing. Record: \texttt{verification/\allowbreak{}lane1/\allowbreak{}refute1/\allowbreak{}VERDICT.md}, scripts r1.py to r5.py.
\begin{enumerate}
\item{} The convex-support case of Caffarelli. Section 2 (1d) now uses only the finite-potential case, plus a weak limit.
\item{} The hopping-difference notation (Step 2): the link is now named.
\item{} Section 6: ``iff'' became ``if''. The ceiling is now called the method's, and the provenance of eps = 0.267 is
given.
\item{} Overclaims: the status line, the monotonicity in L after the 1d table (now shown), and V, which meant both the
volume and the potential (the potential is now Phi).
\end{enumerate}

\item{} Wren's audit (letter 89, 13:01) covered sections 0 to 8 and ran independent spot checks. Sections 0 to 7 HOLD as
written. Section 8 got five fixes, all folded in:
\begin{itemize}
\item{} 8(a) for m $<$ 2 - sqrt2;
\item{} n $<$ 4 and w\_T $<$= 1;
\item{} Adams' eps carried over to H\_T;
\item{} h\_T 4.41 -$>$ 4.42 and the D\_W1/D\_W0 notation;
\item{} c\_f(H\_T) 0.158943 -$>$ 0.158939.
Her polish is also in: the a.e. note in section 0; sigma $>$= $|$lambda\_0$|$, so the statement loses V ln(1 - m\_f\textasciicircum{}2); the
cubed penalty in 1d; and the split line before (3.3).
\end{itemize}

\item{} I checked her numbers before folding them in (sealed \texttt{verification/\allowbreak{}lane1/\allowbreak{}wren09/\allowbreak{}check-plan.md}, log \texttt{out-check.log}); N1 to N5
all hit. One catch of my own: ``mu\_T $>$= 0.3793'' rounded a lower bound up (it is 0.37928), so it now reads 0.3792.
The sealed H\_T check's Q1 used 0.3793 as its threshold, but no sample came near it (minimum 0.3939), so its verdict
stands.
\item{} Round 2 covered Section 8 (and the steps it inherits) and ran 23:49:56 to about 23:56. Its brief held the claim, its use
in Route (v) and the question only. Verdict: PARTIALLY REFUTED. The rate survives. One constant was killed: eps = 0.267
is past eps* = m(2 - m)/(2 + sqrt2) = 0.266533, where Adams' bound reaches zero (-0.00267 at 0.267). The ceilings are
now open bounds, e l $<$ 0.4727 (H\_T) and e l $<$ 0.3367 (H\_W). One wound: L\_0(H\_T) is 32, not 16. Lambda/V $>$ c\_f(H\_T)
until L \textasciitilde{} 4500 is recorded. Survived: c\_f(H\_T), mu\_T(m) for m = 0.1 to 1, the inverse-norm bound under an L-BFGS
ascent (saturated, never exceeded), the H\_T1 gap, h\_T, the n\_0 and T\_0 counts, and c\_f against the free density at
U\_AP (within O(ln L)). Adams' generalisation is still unread first-hand. Record:
\texttt{verification/\allowbreak{}lane1/\allowbreak{}s8-refute/\allowbreak{}} (RULING.md, refuter/VERDICT.md, my\_eps\_check.py). The sweep earlier
today checked the rounding of e l\_0 and never the rounding of the eps feeding it.
\end{itemize}


\end{document}
